%-----------------------------------------------------------------------
%                                                                 aa.tex
% AA vers. 9.3, LaTeX class for Astronomy & Astrophysics
% Demonstration file
%                                                       (c) EDP Sciences
%-----------------------------------------------------------------------
%
%\documentclass[referee]{aa}    % for a referee version
%\documentclass[onecolumn]{aa}  % for a paper on 1 column  
%\documentclass[longauth]{aa}   % for long lists of authors and/or affiliations. 
                                % This command displays the first eight authors on page 1
                                % and shift the whole list after the references.
                                % Ensure to separate each author with the \and command (see below)
%\documentclass[letter]{aa}     % for the letters
%\documentclass[bibyear]{aa}    % if the references are not structured
\documentclass{aa}  
\usepackage{graphicx}
\usepackage{txfonts}
\usepackage{lipsum}
\usepackage{lscape}             % to rotate a single page table, example in appendix.
                                % For landscape tables, see the longtable examples.
\usepackage{placeins}           % useful with \FloatBarrier, to keep 
                                % onecolumn floats from drifting to the next section
\usepackage{natbib}
\usepackage[colorlinks,allcolors=blue,bookmarks=false]{hyperref} 
\usepackage{xcolor}
% Markup for revised text (magenta). For clean submission: \renewcommand{\chg}[1]{#1}
% Note: use \lesssim, \gtrsim throughout (txfonts shortcuts \la, \ga avoided for consistency).
\newcommand{\chg}[1]{\textcolor{magenta}{#1}}
\usepackage{amsmath}
%\usepackage{amssymb}
\hyphenation{get-sources get-filaments get-images get-old non-uniform back-ground}

\newcommand{\draftnote}[1]{\textcolor{magenta}{\textbf{[#1]}}}
\newcommand{\msun}{M_{\odot}}
\newcommand{\Scloud}{\Sigma_{\rm cl}}
\newcommand{\ximax}{\xi_{\rm max}}
\newcommand{\Tbe}{T_{\rm BE}}
\newcommand{\rhoc}{\rho_{\rm c}}
\newcommand{\cs}{c_{\rm s}}
\newcommand{\rzero}{r_{0}}
\newcommand{\Rbe}{R_{\rm BE}}
\newcommand{\rhocloud}{\rho_{\rm cl}}
\newcommand{\Rcloud}{R_{\rm cl}}
\newcommand{\mH}{m_{\rm H}}
\newcommand{\mup}{\mu_{\rm p}}
\newcommand{\muH}{\mu_{\rm H_2}}

% ---------------------------------------------------------------------
% Notation.  Calligraphic letters mark the DERIVED CORRECTION FACTORS --
% the multiplicative quantities the method computes, as distinct from
% anything measured or tabulated.  The family is closed and all-Latin:
% the correction factor itself and the three factors it decomposes into.
% Set \CalFactors to 0 below to typeset them upright instead.
\newcommand{\CM}{\mathcal{C}_{M}}
\newcommand{\Fflux}{\mathcal{F}_{F}}
\newcommand{\Ftemp}{\mathcal{F}_{T}}
\newcommand{\Fbg}{\mathcal{F}_{B}}
% Masses.  \Mex is the mass an extraction reports, \Mbe the true
% Bonnor--Ebert mass of a model, \Mref the known reference mass of a test
% source (\Mbe for a model, the injected mass for an injected core), and
% \Mcorr the mass after correction.
\newcommand{\Mex}{M_{\rm ex}}
\newcommand{\Mbe}{M_{\rm BE}}
\newcommand{\Msed}{M_{\rm SED}}
\newcommand{\Mref}{M_{\rm ref}}
\newcommand{\Mcorr}{M_{\rm corr}}
\newcommand{\Mrec}{M_{\rm rec}}

\begin{document}

\title{Correcting systematic mass biases in prestellar cores\\extracted from \emph{Herschel} images of molecular clouds}
\subtitle{I. Method, physical basis, and validation}

\author{Alexander Men'shchikov\inst{1}}
\institute{D\'epartement d'Astrophysique (DAp), IRFU, CEA Saclay,
Universit\'e Paris-Saclay, 91191 Gif-sur-Yvette, France}

\date{Received / Accepted}
\offprints{Alexander Men'shchikov}
\titlerunning{Correcting systematic mass biases in prestellar cores}
\authorrunning{A.~Men'shchikov}

\abstract{Masses of prestellar cores observed by \emph{Herschel} and derived from fitting their spectral energy distributions are
underestimated through three channels: the dust temperature gradients inside the embedded cores, incorrect determination and
subtraction of background shapes, and inaccurately estimated footprint sizes. These have not previously been combined into a
correction calibrated against a physical core model and validated independently of the calibration sample. We build a grid of
radiative-transfer models of pressure-truncated Bonnor-Ebert (BE) spheres, extract exactly as real cores, and use the resulting
bias to correct real extractions. The correction predicts $\CM=\Mbe/\Mex$ from four observables: background
column, peak-to-mean concentration, footprint-to-size ratio, and the core angular size; the latter breaks a residual degeneracy and
reduces large scatter in $\CM$ among cores sharing the former three observables. The method is validated on three
independent tests: a leave-one-out test on the model grid, an independent model grid spanning the full concentration range, and
model cores injected in and then re-extracted from real \emph{Herschel} images. The four-observable correction is unbiased to
within a few percent, with $85$--$98$\% of sources recovered within a factor of two, and its scatter is at or below the uncorrected
level in every test. The method recovers the true mass of injected, withheld, and leave-one-out cores without tuning to any test.
Application to real clouds, and the resulting revision of core mass function slopes, is presented separately.}

\keywords{Stars: formation -- Infrared: ISM -- Submillimeter: ISM -- Methods: radiative transfer -- Techniques: image processing --
Techniques: photometric}

\maketitle

% =====================================================================
\section{Introduction}
\label{sec:intro}

Far-infrared and submillimeter surveys of molecular clouds with the \emph{Herschel} Space Observatory \citep{Pilbratt_etal2010}
have produced large catalogs of physical properties of prestellar and starless cores \citep[e.g.][]{Andre+2010,Konyves+2015}.
Masses of the cores were obtained by fitting modified-blackbody shapes to the spectral energy distributions (SED) of extracted
cores. The core mass function (CMF) derived from these catalogs is widely compared to the stellar initial mass function (IMF), and
the shape of that comparison depends directly on the accuracy of the individual core masses.

Two physical mechanisms bias the reported SED mass low, and both act in the
same direction. First, a prestellar core is not isothermal: dust heated by
the external interstellar radiation field is warmer at the core's surface
than in its shielded interior, so a single-temperature SED fit is weighted
towards the warm outer layers and under-estimates the mass held by the cold,
optically-thicker interior. Second, the extraction algorithm measures flux
only inside a detection-limited footprint after subtracting a local
background; any core flux outside that footprint, or removed along with the
background, is lost and is not recovered by the SED fit. Both biases grow
with the source's environment -- a denser, more strongly illuminated
background produces both a larger temperature gradient and a smaller
detectable footprint relative to the source size -- so neither is a simple
multiplicative constant.

This paper describes a method to correct for both mechanisms
simultaneously, calibrated against a grid of radiative-transfer models
of embedded cores and validated on synthetic data. Figure~\ref{fig:flowchart}
summarizes the method end to end. Section~\ref{sec:grid} describes the
model grid. Section~\ref{sec:correction} derives the correction algorithm,
in particular the definition of the per-band recovered-flux fraction,
which is the part of the method most sensitive to get right and the part
revised most substantially relative to earlier attempts
(Sect.~\ref{sec:footprint}). Section~\ref{sec:validation} presents the
validation on two independent samples. Section~\ref{sec:discussion}
discusses the method's domain of validity and known limitations.
Section~\ref{sec:conclusions} summarizes. Application of the method to
observed clouds, and its consequence for the core mass function, are the
subject of a companion paper.

\begin{figure*}
  \centering
  \includegraphics[width=0.92\linewidth]{fig_method_flowchart.pdf}
  \caption{Method overview. (\emph{left}, calibration) the radiative-transfer model grid is
  extracted exactly as \emph{getsf} would extract a real source, and the
  self-consistent per-band footprint of Sect.~\ref{sec:footprint}
  (highlighted) converts each truncation into a recovered mass, populating
  a distance-invariant corrector. (\emph{right}, application) three
  distance-invariant observables are measured on a real catalog and
  looked up in the same corrector, for any cloud at any distance. The
  validation results (bottom) are those of Sect.~\ref{sec:validation}.}
  \label{fig:flowchart}
\end{figure*}

% =====================================================================
\section{The model grid}
\label{sec:grid}

\subsection{The Bonnor--Ebert sphere family}
\label{sec:bemodel}

We adopt isothermal Bonnor--Ebert (BE) spheres \citep{Ebert1955,Bonnor1956}
as the core density template. A BE sphere is a self-gravitating,
pressure-supported gas sphere in hydrostatic equilibrium at a single
temperature. Its dimensionless density $e^{-\psi(\xi)}$ obeys the isothermal
Lane--Emden equation,
\begin{equation}
  \frac{1}{\xi^{2}}\frac{{\rm d}}{{\rm d}\xi}
  \!\left(\xi^{2}\frac{{\rm d}\psi}{{\rm d}\xi}\right)
  = e^{-\psi},
  \qquad \psi(0)=0,\ \ \psi'(0)=0,
  \label{eq:lane-emden}
\end{equation}
where $\xi = r/r_0$ is a dimensionless radius. The physical density is
$\rho(\xi) = \rhoc\,e^{-\psi(\xi)}$, with $\rhoc$ the central density and
$r_0$ a characteristic scale length,
\begin{equation}
  \rzero = \frac{\cs}{(4\pi G \rhoc)^{1/2}},
  \qquad
  \cs^{2} = \frac{k_{{\rm B}}\,\Tbe}{\mup\,\mH},
  \label{eq:r0}
\end{equation}
where $\cs$ is the isothermal sound speed at the sphere's temperature
$\Tbe$. Equation~(\ref{eq:lane-emden}) has no closed-form solution and is
integrated numerically once; the resulting function $\psi(\xi)$ is
tabulated and reused for every model in the grid.

A single BE model is fully specified by three numbers: the temperature
$\Tbe$ (through $\cs$), the central density $\rhoc$, and a dimensionless
truncation radius $\ximax$ at which the sphere is cut and joined to the
surrounding cloud. The truncation sets the profile's overall shape and,
through the standard BE stability criterion, whether the truncated sphere
is a stable equilibrium: $\ximax \le 6.451$ is sub-critical (stable),
$\ximax = 6.451$ is the critical point (marginally stable, with a
center-to-edge density contrast of 14.0), and $\ximax > 6.451$ is
super-critical (gravitationally unstable, i.e., collapsing). The
temperature $\Tbe$ and central density $\rhoc$ set only the physical scale
of the profile; $\ximax$ alone carries the shape information. Because
$\Tbe$ and $\rhoc$ do not affect shape, two dimensionless functions of
$\ximax$ suffice to describe the family's projected observables. The first
is the enclosed dimensionless mass,
\begin{equation}
  C(\xi)\equiv\int_{0}^{\xi} e^{-\psi(\xi')}\,\xi'^{2}\,{\rm d}\xi'
  = \xi^{2}\,\psi'(\xi),
  \label{eq:Cxi}
\end{equation}
obtained from the first integral of Eq.~(\ref{eq:lane-emden}) and evaluated
numerically on the tabulated solution. The second is the projected
half-maximum width, $\varphi(\ximax)$, defined as the full width at half
maximum of the line-of-sight column-density profile (the Abel projection
of $e^{-\psi}$) in units of $r_0$, so that a model's intrinsic (unconvolved)
projected width is
\begin{equation}
  H_{\rm BE} = \varphi(\ximax)\,\rzero .
  \label{eq:phidef}
\end{equation}
The function $\varphi(\ximax)$ has a property that matters for the grid's
coverage at high mass: it saturates at large $\ximax$, approaching
$\varphi \simeq 6.6$ as $\ximax \to \infty$, because the projected width is
set by the flat inner core and is nearly insensitive to how far the
diffuse outer envelope extends. A very concentrated sphere is therefore not
much wider, in projection, than a moderately concentrated one at the same
central density; concentration and projected size become progressively
decoupled toward large $\ximax$.

\subsection{Central density and mass from two observables}
\label{sec:inversion}

Each grid model is specified by a target mass $M$ and a target intrinsic
projected width $H_{\rm BE}$ at a chosen temperature
$\Tbe$, rather than by directly choosing $\rhoc$ and $\ximax$. The central
density must be obtained by solving the BE structure for these two
observables, not by renormalizing an arbitrarily scaled profile to a
target mass; renormalization decouples the density from the hydrostatic
solution and produces a profile that is no longer a valid BE sphere.

The enclosed mass of a truncated sphere follows from Eq.~(\ref{eq:Cxi}),
\begin{equation}
  M = 4\pi\,\rhoc\,\rzero^{3}\,C(\ximax) .
  \label{eq:mass}
\end{equation}
Eliminating $\rhoc$ and $\rzero$ between Eq.~(\ref{eq:mass}) and the width
definition, Eq.~(\ref{eq:phidef}), gives a single equation for the
truncation radius that uses only the two observables $M$ and
$H_{\rm BE}$,
\begin{equation}
  \frac{\varphi(\ximax)}{C(\ximax)}
  = \frac{H_{\rm BE}\,\cs^{2}}{M\,G} .
  \label{eq:xiinvert}
\end{equation}
The ratio $\varphi/C$ decreases monotonically with $\ximax$, so
Eq.~(\ref{eq:xiinvert}) has a unique solution for any physically reasonable
$(M,H_{\rm BE})$ pair: a core's mass and size together
fix its concentration. Once $\ximax$ is found (by a numerical root search),
the scale length and central density follow directly,
\begin{equation}
  \rzero = \frac{M\,G}{\cs^{2}\,C(\ximax)},
  \qquad
  \rhoc = \frac{\cs^{2}}{4\pi G\,\rzero^{2}},
  \label{eq:rho_from_MFW}
\end{equation}
and the truncation (outer) radius is $\Rbe = \ximax\,\rzero$. Because
$\varphi$ saturates while $C$ grows without bound as $\ximax$ increases
(Sect.~\ref{sec:bemodel}), a massive core with a small projected size maps
to a very large $\ximax$; this sets the natural high-concentration edge of
the useful grid (Sect.~\ref{sec:gridcoverage}).

\subsection{Why the truncation is a free parameter}
\label{sec:geometry}

An earlier version of this grid fixed $\ximax$ at the critical value,
6.451, on the grounds that only marginally stable spheres are
astrophysically distinguished states. This choice fails observationally.
At fixed temperature, a critical sphere's mass and radius are tied
together through Eqs.~(\ref{eq:mass}) and (\ref{eq:r0}): both scale as
$\rhoc^{-1/2}$, so $M \propto \Rbe$. A more massive critical sphere is
therefore necessarily larger and more diffuse. Compared against a published
core catalog, the fixed-critical family reaches truncation radii of order
$600''$ for its most massive members, while no observed core in the same
catalog exceeds about $110''$; roughly a fifth of the fixed-critical grid
occupies a size regime the observations do not populate. Real massive
cores are compact and centrally concentrated, not large and diffuse, and a
family constrained to $M \propto \Rbe$ cannot represent them.

Freeing $\ximax$ removes this constraint. Solving Eq.~(\ref{eq:xiinvert})
for the truncation that matches each target $(M,H_{\rm BE})$
places every model directly on the size scale it is meant to represent:
higher mass at fixed size is reached by higher concentration (larger
$\ximax$), not by inflating the sphere. The sub-critical, critical, and
super-critical portions of the resulting sequence remain physically
meaningful as templates: the density profile $\rho = \rhoc\,e^{-\psi}$ is a
smooth, monotonically declining function of radius for any $\ximax$, and it
need not be a stable equilibrium to be a reasonable structural template,
because real prestellar cores are not static equilibria either. The
$\ximax$ sequence can be read as tracing an evolutionary progression: weakly
concentrated and often unbound at low $\ximax$, marginally stable near the
critical point, and collapsing above it, where the most centrally
concentrated and unambiguously prestellar cores are expected to lie. The
grid spans this full sequence rather than assuming a single dynamical
state for every core.

\begin{figure}
  \centering
  \includegraphics[width=\linewidth]{fig_grid_coverage.pdf}
  \caption{Coverage of the production grid in the plane of BE mass against
  intrinsic, unconvolved projected size, $H_{\rm BE}$
  (Eq.~\ref{eq:phidef}), given in physical units (parsec, with an
  astronomical-unit axis on top): both are properties of the model itself
  and require no assumed distance. Each point is one of the 518 grid nodes,
  colored by its inverted truncation $\xi_{\max}$ (Sect.~\ref{sec:inversion}),
  obtained by solving Eq.~(\ref{eq:lane-emden}) once and inverting
  Eqs.~(\ref{eq:mass})--(\ref{eq:rho_from_MFW}) for each tabulated
  ($\Mbe$, $\rho_{\rm c}$, $\Tbe$). The red curves are not grid nodes
  but the exact critical sequence, $\xi_{\max}=6.451$ held fixed, computed
  analytically for each of the grid's six background columns $\Scloud$
  (Sect.~\ref{sec:tbe} gives the self-consistency loop). At fixed mass the
  grid spans a range of sizes set by concentration, rather than the single
  mass-size relation, $M\propto\Rbe$, that a family restricted to critical
  spheres alone would trace (Sect.~\ref{sec:geometry}).}
  \label{fig:gridcoverage}
\end{figure}

Figure~\ref{fig:gridcoverage} shows the resulting node coverage in the mass
versus intrinsic-size plane, colored by the inverted $\ximax$ of each node.
Because $\ximax$ is a free parameter, the grid populates a genuine
two-dimensional region of this plane rather than the single curve a
fixed-critical family would trace: at any given mass, nodes span roughly an
order of magnitude in size, from extended, weakly concentrated spheres to
compact, highly concentrated ones.

The red curves mark the critical truncation exactly, rather than
approximating it by the nearest sampled grid nodes. For each of the six
background columns $\Scloud$ used in the grid, we solve, self-consistently
at every mass, for the central density, scale length, and temperature of a
sphere held at $\ximax=6.451$: Eqs.~(\ref{eq:mass}) and~(\ref{eq:r0}) give
$\rho_{\rm c}$ and $r_0$ at a trial $\Tbe$, Eq.~(\ref{eq:Ntot}) gives the
resulting total shielding column, and Eq.~(\ref{eq:TofN}) gives an updated
$\Tbe$; iterating to convergence yields the exact
($\Mbe$, $H_{\rm BE}$) locus of critical spheres at that
$\Scloud$. This construction shows directly why the critical sequence is
not a single line of slope unity, as the naive relation $M\propto\Rbe$ for
a critical sphere at fixed temperature would suggest: the six curves
fan out by about 40\% in size at fixed mass, because $\Tbe$ is not fixed
across the grid but rises with $\Scloud$ (Sect.~\ref{sec:tbe}), and because
$M$ and $\Rbe$ depend on $\Tbe$ through $\cs$
(Eqs.~\ref{eq:mass}--\ref{eq:r0}), eliminating the central density gives
$M\propto\Rbe\,\Tbe$ at fixed $\ximax$, not $M\propto\Rbe$ alone. The
critical sequence cuts diagonally across the full extent of the grid in
both mass and size, and does not by itself bound the useful part of the
grid, consistent with sub-critical, critical, and super-critical models all
serving as valid structural templates (Sect.~\ref{sec:geometry}).

\subsection{Molecular-weight conventions}
\label{sec:mu}

Two distinct mean molecular weights enter the calculation, and using the
wrong one in the wrong place is a common source of error. The mean
molecular weight per free particle, $\mup = 2.33$, enters the sound speed
(Eq.~\ref{eq:r0}) and therefore the entire hydrostatic structure, because
the pressure that supports the sphere is set by the total number density
of free particles, $P = n_{\rm tot}\,k_{{\rm B}}\,T$. The mean
molecular weight per hydrogen molecule, $\muH = 2.8$, enters only where a
mass density or column density is converted to or from a number density of
hydrogen molecules, for example $N_{\rm H_2} = \Sigma/(\muH\,\mH)$.
Using $\muH$ in place of $\mup$ in the sound speed under-estimates $\cs$
and, through Eqs.~(\ref{eq:mass})--(\ref{eq:rho_from_MFW}), the resulting
mass and radius at fixed central density.

\subsection{Radiative transfer and the interstellar radiation field}
\label{sec:rt}

Each BE model is embedded in a uniform surrounding cloud and processed with
\emph{radmc-3d} \citep{Dullemond+2012}, a Monte Carlo dust continuum
radiative-transfer code, to obtain its three-dimensional dust temperature
structure and its emergent surface brightness in the \emph{Herschel} photometric
bands (70, 160, 250, 350, 500\,$\mu$m) and in a high-resolution
($13.5''$) surface-density proxy band -- denoted $\Sigma$, its common
name -- after convolution to each band's point spread function. The
core is heated by the standard interstellar radiation field (ISRF),
tabulated as an interstellar blackbody field plus a scaled hot blackbody
component, normalized to a strength $G_0 \approx 10$ in Habing units;
integrating the far-ultraviolet part of
the field over $4\pi$ steradian and comparing to the Habing normalization
gives $G_0 = 4\pi \times 1.22\times10^{-3}/1.6\times10^{-3} \approx 10$.
Polycyclic aromatic hydrocarbon emission bands are not included, since
adding them changes core-center temperatures by only 1--2\%; the interior
temperature structure is dominated by diffuse dust re-emission, not by
transient PAH heating. $G_0$ is held fixed across the grid, because the
shielded cloud interior that sets a core's heating environment is close to
$G_0$-insensitive at the relevant columns and a per-core $G_0$ is not
recoverable from three-dimensional embedding in any case; the temperature
at which each model is built, $\Tbe$, is instead set self-consistently
(Sect.~\ref{sec:tbe}).

\subsection{The core's construction temperature}
\label{sec:tbe}

The temperature $\Tbe$ that enters the structural equations
(Eqs.~\ref{eq:r0}--\ref{eq:rho_from_MFW}) is the single isothermal
temperature at which the hydrostatic profile is built; it is distinct from
the three-dimensional dust temperature structure $T_{\rm dust}(r)$
that \emph{radmc-3d} subsequently computes for that profile, which is not
isothermal and carries the warm-envelope, cold-center gradient responsible
for the temperature-gradient mass bias this method corrects. Because
$\Tbe$ enters Eq.~(\ref{eq:xiinvert}), it also affects the truncation
$\ximax$ assigned to a given target $(M,H_{\rm BE})$, so it
cannot be chosen arbitrarily; we take $\Tbe$ to be the mass-weighted mean
dust temperature within $\Rbe$, so that each model is built at
(approximately) the temperature it would itself settle to.

The temperature at which a shielded core sits is controlled by the total
radial column between the external radiation field and the core center,
\begin{equation}
  N_{\rm tot} = \Scloud + N_{\rm core},
  \qquad
  N_{\rm core} = \frac{\rhoc\,\rzero}{\muH\,\mH}
                      \int_{0}^{\ximax} e^{-\psi(\xi)}\,{\rm d}\xi ,
  \label{eq:Ntot}
\end{equation}
the sum of the cloud's own column and the core's self-shielding. The
one-sided radial integral, rather than the full line-of-sight column
through the sphere, is the physically appropriate measure because the
incident radiation is isotropic and a grain at the center is shielded along
the radius; it also reproduces the radiative-transfer-computed temperatures markedly
better in practice. Core self-shielding is not a minor correction: across
the grid, $N_{\rm core}$ is typically about 60\% of $\Scloud$ and can
reach $10^{23}\,{\rm cm^{-2}}$ in the most concentrated models, so
$\Tbe$ cannot be tied to $\Scloud$ alone.

We calibrate the relation $\Tbe(N_{\rm tot})$ once, on a dedicated set
of radiative-transfer models spanning the grid's full range of $N_{\rm tot}$, iterated
until the temperature used to build each model agrees with the temperature
\emph{radmc-3d} computes for it. A quadratic in
$x \equiv \log_{10}(N_{\rm tot}/{\rm cm^{-2}})$ reproduces the
calibration set to a root-mean-square of $0.22$\,K,
\begin{equation}
  \Tbe = a_0 + a_1\,x + a_2\,x^2,
  \label{eq:TofN}
\end{equation}
with $(a_0,a_1,a_2) = (1021.4,\,-86.47,\,1.842)$, valid over
$x = 21.5$--$23.4$; the relation runs from $\Tbe \simeq 13.9$\,K at the
lowest columns in the grid to $\simeq 6.5$\,K at the highest. The residual
scatter about this fit is not measurement noise: because $\Tbe$ is a
mass-weighted average over a temperature gradient rather than a function of
the shielding column alone, two models at the same $N_{\rm tot}$ but
different concentration place different fractions of their mass in the
warmer outer layers and so settle to slightly different mass-weighted mean
temperatures. This residual dependence on concentration is small enough
that we retain the single-variable relation of Eq.~(\ref{eq:TofN}).

\subsection{Embedding cloud and dust properties}
\label{sec:cloud}

The surrounding cloud is a uniform-density sphere of radius $\Rcloud$,
large enough that the beam never reaches its boundary and that enough
clean cloud remains around the core for a realistic local background
estimate. Both requirements scale with the core's size as convolved to the
coarsest relevant beam, so we set
\begin{equation}
  \Rcloud = f_{\rm r}\left((O_{500}\,d)^2 + \Rbe^2\right)^{1/2},
  \label{eq:rcloud}
\end{equation}
combining the largest \emph{Herschel} beam and the core's truncation radius in
quadrature, with distance $d$ and $f_{{\rm r}} \approx 30$ chosen so that
the projected cloud column is flat across the core's footprint and its
surrounding background annulus. Here $O_\lambda$ denotes the angular
resolution (beam half-maximum width) of photometric band $\lambda$, with
$O_{160}=13.5''$, $O_{250}=18.2''$, $O_{350}=24.9''$, and
$O_{500}=36.3''$, and $O_\Sigma=13.5''$ for the high-resolution
surface-density proxy band of Sect.~\ref{sec:rt} (numerically equal to
$O_{160}$, but a distinct, composite band), used throughout in place of
the numerical values. The
cloud's line-of-sight column density, rather than its volume density, is
the physically meaningful and model-invariant quantity, since it is what
the observed background measures; we therefore fix the grid's cloud axis to
be the surface density $\Scloud$ (measured just outside the core, to avoid
the core's own contribution) and derive the corresponding uniform volume
density $\rhocloud$ for each model,
\begin{equation}
  \rhocloud = \frac{\Scloud\,\muH\,\mH}{2(\Rcloud^2 - \Rbe^2)^{1/2}} .
  \label{eq:sigmacloud}
\end{equation}
This construction gives identical cloud radiative transfer at a given
$\Scloud$ regardless of core mass or distance, matching the grid axis
directly to what is measured in a real extraction.

For the dust properties, we adopt the \citet{Ossenkopf+Henning1994} opacity
law $\kappa_\nu$, applying a scaling factor of 1.0746 so that it agrees with
the analytic approximation $\kappa_\nu = 10\,(300\,\mu{\rm m}/\lambda)^2\,
{\rm cm^2\,g^{-1}}$, commonly used in \emph{Herschel} Gould Belt Survey SED
fitting, to better than 0.5\% at every \emph{Herschel} band. The same opacity is
used in the radiative-transfer calculation and in the SED mass fit
(Sect.~\ref{sec:footprint}), and the same dust-to-gas mass ratio,
$R = M_{\rm dust}/M_{\rm gas} = 0.01$, is used both to build the
dust density \emph{radmc-3d} requires and to convert a recovered dust mass back to
a total mass; because the correction is inherently differential (a ratio of
two masses computed the same way), the exact value adopted for $R$ cancels
in that ratio and is not a source of bias.

\subsection{Mass-conserving discretization}
\label{sec:massconserv}

\emph{radmc-3d} treats each grid cell as uniform and computes its mass as the
cell's density times its exact shell volume. Evaluating the BE density at
each cell's center, rather than its shell average, does not conserve mass
on a grid coarse enough to be computationally practical; we measured a
mass deficit of about 9\% from center-sampling alone. We instead assign
each cell the density that reproduces its exact enclosed mass, obtained
directly from the function $C(\xi)$ of Eq.~(\ref{eq:Cxi}),
\begin{equation}
  \rho_{\rm cell} = 3\,\rhoc\,
    \frac{C(\xi_{\rm out}) - C(\xi_{\rm in})}
         {\xi_{\rm out}^3 - \xi_{\rm in}^3},
  \label{eq:cellrho}
\end{equation}
with $\xi_{\rm in}$ and $\xi_{\rm out}$ the cell's inner and outer
wall radii. Because adjacent cells share a wall, the enclosed-mass
differences telescope: the total mass summed over all cells equals
$4\pi\,\rhoc\,\rzero^3\,[C(\xi_{\rm edge}) - C(\xi_{\rm inner})]$
exactly, independent of the number of cells, so the grid resolution can be
chosen for photon-noise statistics rather than for mass accuracy. A radial
wall is placed exactly at the truncation radius $\Rbe$, so that no single
cell straddles the BE/cloud boundary and every cell is unambiguously either
pure BE or pure surrounding cloud. With this construction, the integrated
grid mass matches the analytic BE mass of Eq.~(\ref{eq:mass}) to under
$0.1\%$ for every model in the grid (Appendix~\ref{app:massconserv}).

\subsection{Grid coverage and diagnostic flags}
\label{sec:gridcoverage}

The production grid covers cloud backgrounds
$\Scloud=3\times10^{21}$--$9.6\times10^{22}\,{\rm cm^{-2}}$ and BE masses
$\Mbe=0.008$--$27.6\,\msun$, sampled at a nominal factor-of-two node
spacing, refined locally where an observable's interpolation error exceeds
a few per cent (Appendix~\ref{app:interp}). Each node carries two
diagnostic flags recorded for downstream selection: a stability flag read
directly from $\ximax$ (sub-critical, critical, super-critical;
Sect.~\ref{sec:geometry}), and a boundedness proxy
$\alpha_{\rm BE}=M_{\rm BE,crit}/M$ following \citet{Konyves+2015}, with
$M_{\rm BE,crit}=2.4\,R\,c_s^2/G$ the critical BE mass at the model's
radius and sound speed. Neither flag is used to trim the grid: the
extraction detects, and the correction must be able to handle, both bound
and unbound starless cores, as well as collapsing ones.

Models are retained wherever the structure is physical and the resulting
core is observationally detectable; a small fraction of the most compact,
massive models ($\ximax\gtrsim60$) are excluded because their projected size
falls below the beam at every band, making them unresolved point sources
for which no finite truncation radius is observationally distinguishable.
Such cases are flagged as resolution-limited and receive no correction
rather than an unconstrained one.

% =====================================================================
\section{The mass correction}
\label{sec:correction}

\subsection{The observables and the recovery degeneracy}
\label{sec:observables}

The correction predicts $\CM=\Mbe/\Mex$ as a function of
three quantities measurable in any \emph{getsf} extraction, all of which are
either surface brightnesses or dimensionless ratios and therefore carry no
explicit dependence on distance: the local background column density
$\Scloud$ (the $\Sigma$-band peak background, ${\rm PEAK^{BGF}}$), the
peak-to-mean concentration, $\zeta\equiv{\rm PEAK^{SRC}}$ divided by the
footprint-mean column density, and the footprint factor
$\phi=A_{\rm F}/H$, the ratio of the detected source's footprint major
axis, $A_{\rm F}$, to its beam-convolved half-maximum size, $H$
\citep[following the terminology of][]{Menshchikov2021}. Keying the
correction on the reported mass itself, rather than on these three
quantities, was tried and rejected: because reported mass scales
approximately as distance squared, a source observed at different distances
with identical physical structure is matched to different, and increasingly
wrong, grid models as distance grows (Sect.~\ref{app:distance}).

The recovered flux fraction $f_{\rm rec}$ -- the fraction of a source's
total flux that survives inside the detection footprint after background
subtraction -- is not itself an observable, since the true total flux is
precisely the unknown the correction must supply. Left completely
unconstrained, the correction is therefore degenerate: a measured source
of given mass, column, and $\zeta$ is reproduced equally well by a
well-recovered model ($f_{\rm rec}\to1$) and by a poorly-recovered one
truncated close to its peak, and these can imply mass correction factors
differing by an order of magnitude. $\phi$ breaks the worst of this
degeneracy, because it is set by the extraction rather than by the intrinsic
source: a footprint cannot approach the beam, since background fluctuations
broaden any measured intensity distribution, so an anomalously small
$\phi$ cannot occur in practice and is excluded by construction from
the recoverable region of the grid. As Sect.~\ref{sec:fourthobs} shows,
however, $\phi$ does not break it entirely: a materially weaker, but
still consequential, degeneracy remains among sources sharing all three
axes, and removing it is the subject of that section.

\subsection{The per-band footprint model}
\label{sec:footprint}

For each grid node and each $\Sigma$-band truncation radius $r$, a recovered
flux fraction $f_{{{\rm rec}},\lambda}$ is computed for every continuum band
$\lambda\in\{250,350,500\,\mu{{\rm m}}\}$, and the reduced fluxes are refit
to obtain the recovered mass,
\begin{equation}
  \Mrec(r) = M_{\rm SED3bs} \,
  \frac{A(F_\lambda f_{{{\rm rec}},\lambda})}{A(F_\lambda)},
  \label{eq:mrec}
\end{equation}
where $M_{\rm SED3bs}$ is the mass from the untruncated 3-band SED fit,
$F_\lambda$ are the model's true band fluxes, and $A(\cdot)$ is the fitted
modified-blackbody amplitude (proportional to mass) of the 3-band fit to
the argument fluxes. The training sample for the correction is then
$(\Scloud,\,\zeta,\,\phi,\,\CM=\Mbe/\Mrec(r))$,
built by iterating over every node and truncation radius. Getting
$f_{{{\rm rec}},\lambda}$ right is the crux of the method: it is the only
step that converts a purely geometric truncation of the model image into a
mass bias, and two more direct definitions were tried and rejected before
arriving at the one used in production.

$\CM$ can be factored into three physically distinct, independently
measurable pieces, used in the diagnostic of Sect.~\ref{sec:factorscatter}
to attribute where its value comes from: $\CM=\Fflux\times
\Fbg\times \Ftemp$. The flux-recovery factor,
$\Fflux\equiv M_{\rm SED3bs}/\Mrec(r) = A(F_\lambda)/
A(F_\lambda f_{{{\rm rec}},\lambda})$, follows directly from
Eq.~(\ref{eq:mrec}) and captures the mass lost to footprint truncation
alone. The background-subtraction factor,
$\Fbg\equiv F_\Sigma^{\rm crater}/F_\Sigma^{\rm interp}$, is the
ratio of the $\Sigma$-band flux recovered under a background estimated from an
uncontaminated annulus around the source (a ``crater'' background) to the
flux recovered under \emph{getsf}'s standard interpolated background,
which typically over-subtracts under a real source and so gives
$\Fbg>1$; it is measured independently of $\Fflux$, from the
same seven-model background-scan campaign of Sect.~\ref{sec:footprint}.
The temperature-gradient factor, $\Ftemp\equiv \CM/(\Fflux
\Fbg)$, is defined as whatever remains once the other two are
divided out, and isolates the bias from fitting a single temperature to a
source with a real internal temperature gradient.

We derive the band footprint from the physical origin of the enlargement
that \emph{getsf} applies to it. \emph{getsf} measures each band inside a
correspondingly larger footprint than the $\Sigma$-band one, with the median
ratio of the band footprint to the $\Sigma$-band footprint,
$A_{{{\rm F}},\lambda}/A_{{{\rm F}},\Sigma}$, measured at $1.04$, $1.30$, and
$1.90$ for $250$, $350$, and $500\,\mu$m on a seven-model background-scan
injection campaign (189 matched sources). We posit that \emph{getsf}'s footprint in every band is set by
a constant source signal-to-noise ratio (S/N) at the footprint edge; this is
supported empirically, since the median of the background-subtracted source
intensity over the local root-mean-square background fluctuation at the
measured footprint edge is $0.76$, independent of source contrast, in the
same background-scan sample. Writing the dust emissivity per unit column as
$e_\lambda(T)=\kappa_\lambda B_\lambda(T)$, both the source surface
brightness and a structural noise field tied to the same column-density
fluctuations scale as $I_\lambda(s)=N(s)\,e_\lambda(T)$ and
$\sigma_\lambda=\sigma_\Sigma\,e_\lambda(T)$, so that
\begin{equation}
  \frac{I_\lambda(s)}{\sigma_\lambda}
    = \frac{N(s)\,e_\lambda(T)}{\sigma_\Sigma\,e_\lambda(T)}
    = \frac{N(s)}{\sigma_\Sigma},
  \label{eq:sncancel}
\end{equation}
independent of wavelength: the dust SED cancels between source and noise,
and the wavelength dependence of the footprint is set entirely by the
convolution of both fields to each band's beam, not by the dust properties.

For a structural fluctuation field with two-dimensional power spectrum
$P(k)\propto k^\gamma$, the root-mean-square amplitude after convolution
with a Gaussian beam of full width at half maximum $\theta$, restricted to
spatial frequencies above a low-frequency cutoff $k_{\rm low}=1/D$ set by
the footprint diameter $D$, is
\begin{equation}
  \sigma^2(\theta;D) = \int_{1/D}^{\infty} k^{\gamma+1}
    \exp\!\left(-a\,\theta^2 k^2\right){\rm d}k,
  \qquad a=\frac{4\pi^2}{8\ln 2},
  \label{eq:powerspec}
\end{equation}
and the convolution factor is $c_\lambda(D)=\sigma(O_\lambda;
D)/\sigma(O_\Sigma;D)$. \citet{Roy+2019} measure the \emph{Herschel}
$250\,\mu$m power-spectrum slope of Aquila to be $\gamma=-2.26\pm0.1$;
comparable values are $-2.63$ for Polaris and $-3.0$ for a diffuse Taurus
field. We adopt $\gamma=-2.5$; Fig.~\ref{fig:footprint} shows that
$c_\lambda(D)$ varies by under $5\%$ over $-3.0\le\gamma\le-2.0$ at fixed
$D$, so the precise slope is not critical, while the dependence on $D$
itself is strong, particularly for the $500\,\mu$m band once its beam
approaches the footprint scale.

\begin{figure}
  \centering
  \includegraphics[width=0.7\linewidth]{fig_footprint_model.pdf}
  \caption{Convolution factor $c_\lambda(D)=\sigma(O_\lambda;D)/
  \sigma(O_\Sigma;D)$ of Eq.~(\ref{eq:powerspec}), for the
  three continuum bands, as a function of the footprint diameter $D$, at the
  adopted power-spectrum slope $\gamma=-2.5$. Vertical dotted lines mark
  each band's angular resolution, below which the band is unresolved at that
  footprint scale.}
  \label{fig:footprint}
\end{figure}

Combining Eqs.~(\ref{eq:sncancel})--(\ref{eq:powerspec}), and canceling
the (band-independent) edge S/N between the band and $\Sigma$ footprints, the
band footprint radius $R_\lambda$ satisfies, self-consistently,
\begin{equation}
  \frac{I_\lambda(R_\lambda)}{I_\lambda(0)}
    = \frac{\Sigma(r)}{\Sigma(0)}\,
      c_\lambda(2R_\lambda),
  \label{eq:selfconsistent}
\end{equation}
which must be solved for $R_\lambda$ at each truncation radius $r$ and in
each band.

The solution is unique and the equation is well posed. Writing
$g(R)=I_\lambda(R)/I_\lambda(0)-[\Sigma(r)/\Sigma(0)]\,c_\lambda(2R)$, the
first term decreases with $R$ because it is a normalized source profile,
while the second increases with $R$ because a larger footprint admits
fluctuation power from larger scales, so $g$ is strictly decreasing and has
exactly one root. We solve it by bracketed root finding (Brent's method) on
the interval from $r$ to the outermost tabulated radius of that band, which
either converges to the root or reports that none exists inside the
tabulated profile; in the latter case the sample is dropped rather than
assigned a fabricated radius. Over the $27\,441$ solver calls needed to build
the production training set, $93.8\%$ return a bracketed root, $1.8\%$ have no
root within the tabulated profile, and the remainder have no usable profile.

\draftnote{The following paragraph documents a correction to earlier versions of
this work. It should be kept: every number in the paper changed because of
it, and a reader comparing with the earlier preprint needs the explanation.}

Earlier versions solved Eq.~(\ref{eq:selfconsistent}) by fixed-point
iteration, substituting a trial $R_\lambda$ into the right-hand side and
inverting the profile to obtain the next trial. That iteration is not
convergent here. Its map is the composition of a decreasing function with an
increasing one and is therefore decreasing, so wherever the magnitude of its
derivative exceeds unity the iterates oscillate rather than converge, which
is the classical condition for fixed-point divergence. Instrumenting it over
the same $26\,649$ calls, $71.9\%$ oscillated, $23.4\%$ exhausted the
iteration limit without meeting the convergence test and silently returned
the last iterate, and in a further $5.5\%$ the target level fell below the
whole tabulated profile and the radius was clamped to its outer edge. The
consequences were not merely noisy but unphysical: the returned $R_\lambda$
failed to increase monotonically with $r$ in $80$ to $91\%$ of nodes
depending on the band, although enlarging the $\Sigma$-band aperture cannot
shrink the aperture in another band, and the recovered flux fractions
inherited the same defect. The resulting scatter propagated into $\CM$
through the modified-blackbody refit, which converts small errors in the
per-band colors into large errors in the fitted amplitude. Nearest
neighbours in the space of the model's own physical parameters, which by
construction must have nearly identical $\CM$, differed by a factor of
$1.36$; with the bracketed solver they differ by a factor of $1.036$. That
residual factor of $1.036$ is the pipeline's numerical floor, and it sets
the limit against which every scatter quoted in Sect.~\ref{sec:validation}
should be read.

Because $R_\lambda$ sets its own noise cutoff through
$c_\lambda$, Eq.~(\ref{eq:selfconsistent}) automatically enforces
$2R_\lambda>O_\lambda$ for a resolved band without an
externally imposed size cut; a band that remains unresolved after
iteration is dropped from that node's SED refit rather than the node being
discarded. An earlier, fixed exclusion on the $\Sigma$-band footprint diameter
($O_\lambda<D$) discarded $29\%$ of the grid's $\Sigma$-band
truncation samples, all compact cores with $H_\Sigma$ smaller
than the $500\,\mu$m beam; because \emph{getsf} measures each band inside
its own, enlarged footprint rather than the $\Sigma$-band footprint, that
criterion was the wrong one, and Eq.~(\ref{eq:selfconsistent}) recovers the
excluded compact-core population without loss of accuracy
(Sect.~\ref{sec:validation}). The recovered flux fraction entering
Eq.~(\ref{eq:mrec}) is finally the cumulative source integral to
$R_\lambda$,
\begin{equation}
  f_{{{\rm rec}},\lambda} = \frac{\displaystyle\int_0^{R_\lambda}
    I_\lambda(s)\,2\pi s\,{\rm d}s}
    {\displaystyle\int_0^{R_{\rm conv}} I_\lambda(s)\,2\pi s\,
     {\rm d}s}.
  \label{eq:frac}
\end{equation}

\subsection{The estimator}
\label{sec:interp}

The training samples of Sect.~\ref{sec:footprint} give the value of $\CM$ at
a scattered set of points in the space of the observables. By a training
sample we mean one entry in the set from which the corrector is built: one
model of the grid evaluated at one truncation radius, giving one value of
each observable and one value of $\CM$. The training set is the collection
of all of them, 7497 samples drawn from 403 nodes of the model grid. To correct a real
source we need the value of $\CM$ at a point where no training sample lies,
so we need a rule that turns the scattered samples into a function defined
everywhere. Two such rules were tested, and the choice between them matters
more than one might expect.

The first rule is piecewise linear interpolation on a Delaunay
triangulation. The training points are joined into a mesh of simplices --
triangles in two dimensions, tetrahedra in three, and their
four-dimensional analogue here -- and the value at a query point is the
linear interpolation between the vertices of the one simplex that contains
it. This is the natural choice for scattered data and it was used in earlier
versions of this work. It has two drawbacks. It uses only the handful of
vertices of a single simplex, so it inherits the error of those few samples
rather than averaging over many; and it is undefined outside the convex hull
of the training set, returning no value at all for a source whose
observables fall beyond the region the grid covers.

The second rule is local linear regression. To evaluate the correction at a
query point, the $K$ nearest training samples are found in a space where
each axis has been divided by its own standard deviation, so that the
distance does not depend on the units of the axes; a plane is then fitted to
those $K$ samples by weighted least squares and evaluated at the query
point. The weights follow the tricube function
$w=\left[1-(d/h)^{3}\right]^{3}$, where $d$ is the distance to a training
sample and $h$ the distance to the furthest of the $K$, so that nearby
samples count most and the weight falls smoothly to zero at the edge of the
neighborhood. Averaging over many samples suppresses the sample-to-sample
noise that the simplex interpolation transmits directly, and because the
fitted plane can be evaluated anywhere, a query outside the hull receives a
value rather than none.

The two rules are compared by leave-one-out validation, which tests a
corrector on data it has not seen. One node is removed from the training set
together with all of its samples, the corrector is rebuilt from what
remains, and it is then asked to predict $\CM$ for the removed node;
repeating this for every node gives one honest prediction per node. Removing
the whole node rather than a single sample is essential, because samples
from the same node differ only in truncation radius and are therefore not
independent of one another, so leaving out a single sample would leave
near-duplicates of it in the training set and the test would measure almost
nothing. Measured this way (Sect.~\ref{sec:validation}), the local linear
estimator reduces the root-mean-square residual in $\log_{10}\CM$
from $0.0583$~dex (a factor of $1.144$) to $0.0445$~dex (a factor of
$1.108$), and returns a value for all $403$ nodes rather than for the $388$
that fall inside the hull. The improvement is insensitive to the
neighborhood size over the range $K=80$ to $640$, so $K$ is not a tuned
parameter; we adopt $K=320$. $K$ is a hyperparameter, that is a setting of
the estimator itself rather than a quantity fitted from the data, and
choosing one by inspecting the same test used to report the final
performance would flatter the result. We therefore also verified it by
split-half selection, in which the choice is made on a randomly selected
half of the nodes and its performance reported on the other half, which the
choice never saw; that gives $0.0437$~dex, confirming the improvement is not
an artifact of the choice. Rescaling the axes before triangulating,
without changing the estimator, gives $0.0582$~dex, that is no improvement
at all, which shows that the deficiency of the Delaunay rule is not its
metric but its reliance on a single simplex.

A plane fitted to $K$ samples can be evaluated arbitrarily far from them,
and the value it returns is then an extrapolation of a first-order fit
rather than an interpolation. We therefore bound each prediction by the
range of $\CM$ spanned by the training samples the fit was built from. On
the model grid this bound never acts, and the leave-one-out residual is
unchanged to four decimal places. On real catalogs it acts for a few per
cent of sources in most clouds, and it is what keeps the correction of a
source unlike anything in the grid finite rather than unbounded.

Earlier versions of this work applied two physical bounds to the output:
$\CM\ge1$, on the argument that every mechanism corrected for adds mass and
never removes it, and a cap on the flux-recovery factor at
$1/{\rm MIN\_FRAC}$ with ${\rm MIN\_FRAC}=0.25$. Both have been
removed. With the footprint solver of Sect.~\ref{sec:footprint} corrected,
they change the leave-one-out scatter from a factor of $1.145$ to $1.15$ and
the median corrected mass ratio from $1.007$ to $1.008$, so they no longer
do anything; they were compensating for a numerical defect rather than
enforcing physics. The argument behind the first is also not supported by
the training data, in which $7.2\%$ of nodes have a true $\CM$ below unity.
\draftnote{Consequence for catalogs: individual sources may now receive $\CM<1$.
This is the honest behavior but it will look different from earlier
catalogs.}

${\rm MIN\_FRAC}$ retains a second, unrelated role earlier in the
pipeline, at the per-band level rather than the output level. The 3-band
mass refit of Eq.~(\ref{eq:mrec}) is a least-squares fit over a temperature
grid, and it is not automatically robust: when the per-band recovered
fractions $f_{{{\rm rec}},\lambda}$ of the bands entering the fit are highly
imbalanced -- one band recovering under $20\%$ of its flux while another
recovers essentially all of it -- the reduced-flux vector can take a shape
no physical modified blackbody can reproduce, and the least-squares solution
is driven to an unphysical temperature at the edge of the fitted range, with
an amplitude, and hence a recovered mass, wrong by up to three orders of
magnitude. The same instability, in a different form, affects a subset of
nodes where only two bands remain resolved: two data points exactly
determine two fit parameters, temperature and amplitude, with no redundancy
to check the result against. We require at least three bands with
$f_{{{\rm rec}},\lambda}\ge{\rm MIN\_FRAC}$ for a sample to enter the
training set at all; a truncation radius failing this requirement is
dropped, and a node all of whose truncation radii fail it is absent from the
grid entirely. This removed $92\%$ of the catastrophic ($>\!10\times$) mass
errors identified in the production grid, at a cost of $7\%$ of the training
samples, and it is applied uniformly, before either the three- or
four-observable corrector of Sect.~\ref{sec:fourthobs} is built.

\subsection{A fourth observable: angular size}
\label{sec:fourthobs}

$\phi$ (Sect.~\ref{sec:observables}) rules out the most extreme degeneracy,
in which a source's true flux is essentially unconstrained, but it does not
determine $\CM$. We tested this directly. For every training sample we
located its nearest neighbor, in the normalized
$(\Scloud,\,\zeta,\,\phi)$ space, drawn from a different physical
node, so that the comparison is between genuinely different cores rather
than between two truncations of the same one. If the three observables
determined $\CM$, then the difference in true $\CM$ between increasingly
close neighbors would go to zero. Instead, binning pairs by separation, the
root-mean-square difference in $\log_{10}\CM$ is flat at $0.14$ to
$0.17$~dex from the closest decile of pairs to the farthest
(Fig.~\ref{fig:degeneracy}). That is the signature of a degeneracy rather
than of interpolation error. The two are distinct and the distinction
matters throughout what follows. A degeneracy means that the map from the
observables to $\CM$ is not one-to-one, so that two sources with the same
observables can have genuinely different correction factors and no rule
whatsoever can tell them apart; interpolation error is the error made
because the training samples are finite in number and the rule must guess
between them. Interpolation error shrinks as the samples used get closer to
the query, whereas a degeneracy does not, which is exactly what the test
above measures. Quantitatively, the three observables account for about
$15\%$ of the variance of $\log_{10}\CM$ across the training set, the
variance explained being one minus the ratio of the residual variance to the
total variance; the remaining $85\%$ is invisible to them. A corrector that
returned the population median for every source would explain no variance at
all, however close its median happened to be to the truth.

The missing information is the source's absolute scale. The three existing
observables are all ratios or surface brightnesses, constructed to be
scale-free (Sect.~\ref{sec:observables}); two models with the same shape but
different size are, by design, indistinguishable to them. But $\CM$ is not
scale-free: it depends on the temperature structure and flux recovery that
follow from the model's actual size, through the self-shielding column
$N_{\rm core}\propto\rho_{\rm c}\,r_0$ of Eq.~(\ref{eq:Ntot}). We therefore
add a fourth axis carrying scale.

\begin{figure}
  \centering
  \includegraphics[width=\linewidth]{fig_degeneracy.pdf}
  \caption{Root-mean-square difference in the true correction factor $\CM$
  between pairs of training samples drawn from different physical models, as
  a function of how close the pair is in the observables (normalized,
  dimensionless units), binned by decile of pair separation. For the three
  observables $(\Scloud,\,\zeta,\,\phi)$ the difference does not shrink as
  the pair separation goes to zero, showing that they do not, even in the
  limit of perfect sampling, determine $\CM$. Adding the angular size as a
  fourth axis halves the difference from the farthest to the closest decile,
  which is the behavior of ordinary interpolation error rather than of a
  degeneracy.}
  \label{fig:degeneracy}
\end{figure}

\subsubsection{The axis is an angle, not a physical length}

The scale-bearing quantity available on the model side is the truncation
radius $\Rbe$ (Sect.~\ref{sec:bemodel}); on the real-source side it is the
footprint radius $(A_{\rm F}B_{\rm F})^{1/2}/2$. Both are measured
angles. Whether the axis should be used as an angle, or converted to a
physical length using the source's distance, is not a matter of convenience:
the two choices give different answers for any cloud at a distance other
than that at which the grid was built, and only one of them is correct.

The correct choice follows from what $\CM$ is. The mass of a source is
$M\propto F_\nu d^{2}/[\kappa B_\nu(T)]$, and the same factor $d^{2}$
appears in the true mass and in the reported mass, so it cancels exactly in
their ratio. Surface brightness does not depend on distance either, and both
the beam and the structural noise that fixes the footprint
(Sect.~\ref{sec:footprint}) are properties of the map in its angular frame.
Every quantity on which $\CM$ depends is therefore an angle, and $\CM$
cannot depend on how far away the source lies. We accordingly use the axis
in angular units on both sides, expressing the model's $\Rbe$ as the angle
it subtends at the distance at which the grid was rendered, and using the
measured footprint radius of a real source directly.

This is confirmed by the data, in a test that belongs to the companion
paper because it uses observed clouds rather than models. Applying the
corrector to seven clouds spanning $139$ to $2000$~pc, the median correction
factor ranges over a factor of $1.70$ and correlates with $\log_{10}d$ at
$0.958$ when the axis is converted to parsecs at each source's own distance,
against a factor of $1.12$ with a correlation of $0.464$ when it is used as
an angle. The trend in the first case is an artifact of the conversion,
because the median angular footprint radius is nearly the same in every
cloud, $28$ to $44$~arcsec, being set by the extraction and the beam rather
than by the sources, so the converted physical radius scales almost purely
with distance. Measured at fixed angular footprint radius, the concentration
$\zeta$, the footprint factor $\phi$, and the source contrast show no
ordering with distance at all.

Because the whole grid was rendered at a single distance, the model-side
axis is the same variable whether expressed in parsecs or in arcsec, and no
test performed on the grid alone can distinguish the two. The distinction
appears only when a real source is looked up. This is why the point was
missed in earlier versions of this work, which used the physical form.

One simplification follows. The correction factor no longer requires knowing
the cloud's distance at all. Distance still sets the mass scale, since the
reported mass carries $d^{2}$, but it no longer enters the correction, so a
distance error propagates as a clean multiplicative factor rather than also
displacing the source along a corrector axis.

\subsubsection{Performance and alternatives}

With the angular size as the fourth axis, the difference in true $\CM$
between nearest neighbors from different nodes falls from $0.087$~dex in the
farthest decile of pair separation to $0.045$~dex in the closest, a factor
of two, which is the behavior of interpolation error and not of a
degeneracy. Subtracting the pipeline's own numerical floor in quadrature,
the irreducible degeneracy of the four observables is about $0.030$~dex, a
factor of $1.07$, against a total spread of $\log_{10}\CM$ of $0.142$~dex.
About $95\%$ of the variance of $\log_{10}\CM$ is therefore in principle
reachable with these four observables, against $15\%$ with three.

We compared several alternative fourth axes, all measurable on both the
model side and in a getsf catalog, under the same leave-one-out test and
with each used consistently for training and for the query. In order of the
root-mean-square residual in $\log_{10}\CM$ they are: a Bonnor--Ebert
profile fit to the source's own radial profile ($0.086$~dex), the reported
mass ($0.098$), the footprint radius ($0.118$), the deconvolved half-maximum
size ($0.118$), the half-maximum size ($0.119$), the SED temperature
($0.125$), the ratio of the $500$ to the $250\,\mu$m half-maximum sizes
($0.126$), no fourth axis at all ($0.140$), and the peak column density
($0.142$). The last of these is worse than using no fourth axis, and the two
temperature-like quantities help only marginally.

We adopt the footprint radius. It is not the best performer of that list,
but it requires nothing beyond the standard catalog, whereas the profile fit
requires the per-source radial profile files, which not every extraction
produces. For extractions that do have them, that axis is obtained by fitting
the analytic Bonnor--Ebert profile of Sect.~\ref{sec:bemodel}, projected and
convolved to the band resolution, to the source's own azimuthally averaged
radial profile by weighted least squares, with the truncation $\ximax$ and the
scale radius $\rzero$ free and the profile taken as \emph{getsf} writes it,
without rim subtraction. Its accuracy is governed by how far the fitted range
reaches relative to the true $\Rbe$: on the model grid the ratio of fitted to
true $\Rbe$ has a $16$th-to-$84$th percentile range of $0.949$--$1.359$ where
the fit extends to $0.7$--$1.0\,\Rbe$, tightening to $0.904$--$1.039$ beyond
$1.5\,\Rbe$ and degrading badly short of $0.4\,\Rbe$. The non-deblended
surface-density profile is therefore the one to use, reaching a median of
$105''$ against $27''$ for the deblended profile of the same sources.
\draftnote{Implementation not preserved beyond the exploratory stage; the
$0.086$~dex figure above comes from that exploration. Either restore the
script before submission or present the number as indicative.}

% =====================================================================
\section{Validation}
\label{sec:validation}

The correction is validated on three independent tests, none used to
calibrate the grid, the footprint model of Sect.~\ref{sec:footprint}, or
the fourth observable of Sect.~\ref{sec:fourthobs}. Each test is run for
both the three-observable corrector and the four-observable corrector, so
that the benefit of the fourth axis is measured on the same samples rather
than inferred indirectly. Two of the three tests are model-grid tests, in
which a node (one physical radiative-transfer model, at one $\Scloud$, mass, and
$\ximax$) is evaluated at many different $\Sigma$-band truncation radii $r$
(Sect.~\ref{sec:footprint}), each giving a distinct
$(\Scloud,\zeta,\phi)$ combination and hence a
distinct test case; we call each such (node, truncation radius)
combination a sample, to distinguish it from the physical node it
comes from. A single node therefore contributes several samples (a median
of $20$ per node in the production grid), not one; where this matters for
an unbiased sample count, we also report results reduced to one
representative sample per node (the sample whose $\phi$ is closest
to $1.84$, the grid's own median value).

Population injection test. Synthetic BE cores are injected into real Aquila
maps and
re-extracted with \emph{getsf}, exactly as a real source would be, and the
resulting catalog is passed through the same correction pipeline used for
real data, using the known distance to Aquila ($260\,$pc) for the fourth
observable. This is the most realistic test available, since it exercises
the full observational chain -- injection, extraction, background
subtraction, SED fitting -- rather than only the model grid. Each of the
108 sources is a single, independent injection, with no node/sample
distinction. How the injected population is chosen, and why that choice
determines what the test is able to measure, is set out in
Sect.~\ref{sec:injpop}.

Independent subdivision grid. A second radiative-transfer grid, disjoint from the
production grid, is built at five background columns interleaved between
the six production-grid values (the geometric mean of each adjacent pair),
each populated with the same target mass grid used in the production grid
but spanning the full range of truncation $\ximax$ from sub-critical
through critical to strongly super-critical, rather than a single
stability class. This gives 226 detectable models (of 314 attempted; the
remainder fall below the detectability floor and are not part of the
test), $3416$ truncation-radius samples in total, or $166$ reduced to one
sample per node. This is a deliberate improvement over an earlier version
of this test, which used $36$ models confined to the critical truncation
only ($\ximax=6.451$ throughout) and consequently exercised only a
one-dimensional slice of the grid's concentration range; the present
version spans the same stability diversity as the production grid itself.
This test isolates the accuracy of the correction surface, independent of
any extraction or injection systematics.

Node-level leave-one-out. For each of the $406$ nodes of the
production grid that retain at least one sample after the reliability
requirement of Sect.~\ref{sec:interp}, every sample belonging to that node
(all of its truncation radii together) is removed, both correctors are
rebuilt from the remaining samples, and the correction is evaluated on the
excluded node's own representative sample (again the one closest to
$\phi=1.84$), giving $388$ independent, genuinely unseen test
cases. This differs from the subdivision test in scale and in what it
stresses: it uses the same physical grid rather than a second,
independently generated one, but it repeats the withhold-and-predict
procedure at (in effect) every node rather than at a fixed external
sample, so it is far less sensitive to whether that external sample
happened to sit in an easy or a hard part of the grid. We restrict \emph{loo} to
one sample per node throughout, rather than reporting all truncation
radii of the withheld node as in the other two tests: the several
truncation-radius samples of one withheld node are correlated (a
prediction that is biased for that node tends to be biased for all of
them), so counting them as independent test cases would overstate the
number of genuinely independent checks the test performs.

Background-scan injection test. A complementary, fourth test asks a
different question from the other three: not whether the correction
generalizes across the grid's diversity of physical models, but whether it
handles the variation a single fixed model can show under different real
background conditions. Seven specific production-grid nodes are each
injected at $\sim\!40$ different positions across the real Aquila field
and re-extracted with \emph{getsf}, giving $85$ matched sources (of $256$
attempted; most losses are sources failing \emph{getsf}'s own
QUALITY$=$ok criterion, not a footprint-radius or model-hull failure). By
construction this test cannot exercise the observable degeneracy of
Sect.~\ref{sec:fourthobs} (removing it requires distinguishing
different models sharing the same three observables, and here the
underlying model is always one of only seven, held fixed), so it is not
expected to show the same pattern as the other three; we report it
separately rather than pooling it with them.

\begin{figure*}
  \centering
  \includegraphics[width=0.72\linewidth]{fig_validation.pdf}
  \caption{Validation of the three- versus four-observable correction, on
  four independent tests. (\emph{top to bottom}) injection--recovery
  ($n=109$, diverse models); independent subdivision grid ($n=3416$
  truncation samples, panel d reduced to $166$ per-node samples for the
  CDF); node-level leave-one-out ($n=388$, one sample per node);
  background-scan injection ($n=85$, seven fixed models at varying real
  backgrounds). (\emph{left}) corrected mass ratio for the three-observable
  (open gray) and four-observable (filled red) corrector versus the
  reference mass; the ideal ratio of unity is dashed. (\emph{right}) the
  empirical cumulative distribution of $|\log_{10}({\rm ratio})|$
  for each corrector; the vertical dotted line marks a factor of two.}
  \label{fig:validation}
\end{figure*}

Table~\ref{tab:validation} and Fig.~\ref{fig:validation} summarize the
result. In every test, both correctors move the median ratio from a
significant low bias to within a few per cent of unity: systematic-bias
removal is not specific to the fourth observable, and both correctors
achieve it. What differs is everything else. The three-observable
corrector's accuracy is mixed: it improves the fraction within a factor of
two over the uncorrected baseline in the subdivision test ($87\%\to90\%$)
and the background-scan test ($79\%\to94\%$), but not in the injection
test ($85\%\to83\%$) or \emph{loo} ($87\%\to88\%$, essentially unchanged), and its
scatter (expressed as a multiplicative factor) is worse than uncorrected
in every test except background-scan (injection $2.06\to2.25$; subdivision
$1.56\to1.58$; \emph{loo} $1.51\to1.62$; background-scan $1.43\to1.37$). The
four-observable corrector's fraction within a factor of two improves over
uncorrected in three of four tests ($85\%\to85\%$ injection,
$87\%\to98\%$ subdivision, $87\%\to95\%$ \emph{loo}, $79\%\to87\%$
background-scan), and its scatter is better than uncorrected in three of
four and closely comparable in the fourth (injection $2.06\to2.09$;
subdivision $1.56\to1.32$; \emph{loo} $1.51\to1.42$; background-scan
$1.43\to1.41$). Compared directly against the three-observable corrector
on the same samples, the fourth observable reduces the typical error
(median $|\log_{10}\text{ratio}|$) in every test -- by roughly a fifth in
the injection test ($0.101\to0.082$) and the background-scan test
($0.107\to0.084$), and by a factor of two or more in the two model-space
tests ($0.042\to0.019$ subdivision; $0.052\to0.031$ \emph{loo}) -- and improves
the individual prediction for the majority of sources in three of four
tests (injection $60\%$; subdivision $74\%$; \emph{loo} $66\%$; background-scan
$59\%$).

The background-scan test's pattern is visibly different from the other
three (Fig.~\ref{fig:validation}g, h): both correctors leave a substantial
residual bias (median $0.82$--$0.84$, versus $0.96$--$1.02$ elsewhere),
and the three-observable corrector is competitive with, or on some metrics
better than, the four-observable one. This is consistent with what the
test can and cannot probe: with only seven distinct underlying models,
repeated many times, there is little room for the specific failure mode
Sect.~\ref{sec:fourthobs} identifies (different models sharing the same
three observables) to occur, so the fourth observable's main advantage is
largely absent by construction here. The residual bias common to both
correctors in this test is a separate finding, not yet explained, and is
noted as a target for follow-up rather than folded into the headline
comparison.

\begin{table}
\caption{Validation of the three- and four-observable correctors on the four
independent tests of Sect.~\ref{sec:validation}. Tabulated, for each test and
each treatment, are the median of the ratio of the mass obtained to the
reference mass, the scatter of that ratio expressed as a multiplicative
factor, and the percentage of sources whose ratio lies between $0.5$ and $2$.
All entries are restricted to sources inside the domain of the correction, that
is to those for which both corrections are interpolations between grid models
rather than extrapolations beyond them (Sect.~\ref{sec:domain}). The
subdivision and \emph{loo} tests use one representative sample per node.}
\label{tab:validation}
\centering
\begin{tabular}{lcccc}
\hline\hline
test & quantity & uncorr. & 3-axis & 4-axis \\
\hline
Node-\emph{loo}      & median           & $0.84$ & $1.13$ & $1.01$ \\
($n=388$)            & scatter (factor) & $1.36$ & $1.35$ & $1.11$ \\
                     & within $2\times$ & $89\%$ & $95\%$ & $100\%$ \\
\hline
Subdivision          & median           & $0.84$ & $1.17$ & $1.02$ \\
($n=156$)            & scatter (factor) & $1.33$ & $1.34$ & $1.11$ \\
                     & within $2\times$ & $90\%$ & $96\%$ & $100\%$ \\
\hline
Bg-scan              & median           & $0.76$ & $1.09$ & $0.93$ \\
($n=110$)            & scatter (factor) & $1.98$ & $1.94$ & $1.95$ \\
                     & within $2\times$ & $79\%$ & $88\%$ & $85\%$ \\
\hline
Population injection & median           & $0.79$ & $1.13$ & $0.93$ \\
($n=95$)             & scatter (factor) & $1.56$ & $1.60$ & $1.56$ \\
                     & within $2\times$ & $85\%$ & $84\%$ & $88\%$ \\
\hline
\end{tabular}
\tablefoot{The subdivision and \emph{loo} ratios are $M/\Mbe$; the bg-scan and
population-injection ratios are $M/\Mref$, the mass of the model injected. The
ideal median is unity. The population-injection entries are for the
flattened, isolated variant, in which the cloud fluctuations beneath each core
and the neighbouring injected cores are both absent, pooled over four
realizations; the effect of restoring each is given in
Table~\ref{tab:contamination}.}
\end{table}

\subsection{The domain of the correction}
\label{sec:domain}

The correctors interpolate a grid of radiative-transfer models. A source whose
observables lie inside the region the models span is corrected by
interpolation between them; one that lies outside is corrected by extrapolation
of a fitted plane, which is a far weaker statement. The two must be
distinguished, and the pipeline reports the distinction per source: a convex
hull is built once from the training samples in each set of observables, and a
source is flagged as inside the domain when both the three- and the
four-observable queries fall within it.

The flag matters in practice. Applying the corrector to seven clouds spanning
$130$ to $1700$~pc, $3820$ of their $8106$ sources, that is $47\%$, lie outside
the hull. They do
so for one reason: $3287$ of them fall below the lowest embedding column of
the grid, $3\times10^{21}\,{\rm cm^{-2}}$, while no more than $99$ fail on any
other axis. The limitation is thus a gap in the coverage of the present grid
at faint backgrounds rather than a property of the method, and it is closed by
extending the grid downward in embedding column.

We report the corrected masses of the extrapolated sources, since the
extrapolation is bounded by the range of the neighbouring training samples and
does not diverge, but the flag is carried in the output catalogue so that a
user may restrict the sample at will. All validation figures quoted in this
paper are for sources inside the domain.

\subsection{Fitting the spectral energy distribution: how many bands}
\label{sec:bands}

A modified blackbody has two free parameters, so a fit to two photometric
points has no degrees of freedom. It passes exactly through both points however
badly they are measured, its chi-squared is zero by construction, and no
goodness-of-fit criterion can reject it. The injections make the consequence
measurable, because the true mass of every recovered core is known.

Table~\ref{tab:bands} gives the outcome as a function of the two settings that
govern which bands are fitted: the factor by which a band's peak and total
fluxes must exceed their own uncertainties before that band is admitted, and
the number of admitted bands below which no fit is attempted. Two conclusions
follow. Fits to two bands are unsafe at every threshold, failing by more than
a factor of five for between $3\%$ and $13\%$ of cores and never reaching zero;
raising the threshold does not repair them, because the defect is the absence
of degrees of freedom rather than the quality of the fluxes. And a threshold of
unity, which is a natural choice and has been used in this context before, is
not sufficient even with three bands, leaving $6\%$ of cores wrong by more than
a factor of five.

We therefore admit a band when its peak and total fluxes both exceed $1.5$
times their uncertainties, and fit only when at least three bands are admitted.
On the injected cores this setting leaves one core in a hundred wrong by more
than a factor of five, against six in a hundred at a threshold of unity, while
retaining $230$ cores against the $135$ that survive a threshold of two.

One feature of Table~\ref{tab:bands} deserves emphasis. At the loosest setting
the median ratio is $1.003$, which appears to be an excellent result, yet the
scatter is a factor of $3.87$ and one core in eight is wrong by more than a
factor of five. The apparent accuracy is the cancellation of large errors of
both signs. The genuine underestimate that this paper exists to correct, near
$0.8$, becomes visible only once the failed fits are removed.

\begin{table}
\caption{Accuracy of the mass reported by the extraction, as a function of the
two settings that govern which bands enter the fit of the spectral energy
distribution. A band is admitted only when its peak flux and its total flux
each exceed their own uncertainty by the factor in the first column; a fit is
attempted only when at least the number of bands in the second column are
admitted. Tabulated are the number of injected cores that survive the setting,
the median of $\Mex/\Mbe$ over them, its scatter as a multiplicative factor,
the percentage within a factor of two, and the number and percentage whose
ratio lies outside $0.2$ to $5$. Measured on the four population-matched
injection realizations, in which the true mass of every recovered core is
known. The adopted setting is in bold.}
\label{tab:bands}
\centering
\begin{tabular}{ccrccc}
\hline\hline
S/N & bands & cores & median & scatter & wrong by $>\!5\times$ \\
\hline
$1.0$ & $2$          & $427$ & $1.003$ & $3.87$ & $56$ ($13\%$) \\
$1.0$ & $3$          & $331$ & $0.906$ & $2.30$ & $20$ ($6\%$) \\
$1.0$ & $4$          & $195$ & $0.792$ & $1.62$ & $0$ \\
$1.5$ & $2$          & $329$ & $0.885$ & $2.95$ & $30$ ($9\%$) \\
$\mathbf{1.5}$ & $\mathbf{3}$ & $\mathbf{230}$ & $\mathbf{0.815}$
      & $\mathbf{1.77}$ & $\mathbf{2}$ ($\mathbf{1\%}$) \\
$1.5$ & $4$          & $108$ & $0.751$ & $1.46$ & $0$ \\
$2.0$ & $2$          & $228$ & $0.817$ & $2.24$ & $9$ ($4\%$) \\
$2.0$ & $3$          & $135$ & $0.751$ & $1.52$ & $0$ \\
$2.0$ & $4$          & $55$  & $0.726$ & $1.31$ & $0$ \\
$3.0$ & $3$          & $63$  & $0.729$ & $1.49$ & $0$ \\
\hline
\end{tabular}
\end{table}

\subsection{Cloud fluctuations and blending}
\label{sec:contamination}

Two contaminations act on an extracted core and neither is a property of the
core itself: the fluctuations of the cloud beneath it, and the light of a
neighbouring source. Their contributions can be separated because the injected
population is built once and then rendered four times. Every core is injected
at the same position with the same model in all four; the runs differ only in
whether the background under each core is flattened before insertion, and
whether the cores that overlap a neighbour are removed. Any core surviving both
filters therefore appears in all four fields, and the comparison is made core
by core rather than sample against sample.

That distinction is not pedantic. Compared sample against sample, an earlier
pair of runs appeared to show the cloud fluctuations biasing the recovered mass
high by $39\%$; the effect vanished when the same cores were compared in both,
the apparent bias having been the different sets of cores that each run
happened to recover.

Table~\ref{tab:contamination} gives the result over four realizations. The
cloud fluctuations inflate the scatter by a factor of about $1.2$ to $1.3$ and
shift the median by an amount consistent with zero. Blending shifts the
recovered mass downward by about $2\%$, small but established at a probability
of $8\times10^{-4}$, and the two do not simply add: the blending shift is
weaker when the fluctuations are also present, so the errors they introduce are
not independent.

\begin{table}
\caption{The two contaminations, measured core by core. Each row compares two
of the four injection runs on the cores recovered in both, the runs being
identical except in the single respect named. The scatter columns give the
standard deviation of $\log_{10}(\Mex/\Mbe)$ in each run, expressed as a
multiplicative factor; the shift is the median of the paired difference of that
logarithm, first run minus second, and $p$ is the probability from a Wilcoxon
signed-rank test that the paired differences are centred on zero. Pooled over
four realizations of $600$ injected cores each.}
\label{tab:contamination}
\centering
\begin{tabular}{lrcccc}
\hline\hline
comparison & pairs & scatter 1 & scatter 2 & shift & $p$ \\
\hline
\multicolumn{6}{l}{\emph{Cloud fluctuations present versus removed}} \\
neighbours present & $113$ & $1.84$ & $1.58$ & $-0.006$ & $0.99$ \\
neighbours removed & $76$  & $1.96$ & $1.46$ & $+0.029$ & $0.07$ \\
\hline
\multicolumn{6}{l}{\emph{Neighbouring cores present versus removed}} \\
fluctuations present & $108$ & $2.34$ & $2.40$ & $-0.001$ & $0.02$ \\
fluctuations removed & $95$  & $1.57$ & $1.56$ & $-0.009$ & $0.0008$ \\
\hline
\end{tabular}
\tablefoot{The shift is in dex. A negative shift means the first run reports
less mass than the second for the same core.}
\end{table}


\subsection{Constructing the injected population}
\label{sec:injpop}

The injection--recovery test is the only one of the four that exercises the
full observational chain, and it is therefore the one whose sample must be
built with most care: the cores that are injected determine which part of the
method's domain the test actually probes. Four requirements govern that
construction. They follow from the structure of the problem rather than from
anything specific to Aquila, so we state them here and give the procedure that
implements them in Appendix~\ref{app:injection}.

\subsubsection{Selection must not use the quantity under test}
The natural way to make
an injected sample resemble a real one is to match the injected cores to real
cores by mass. But the reported mass is what the method is being tested for,
and it is biased; that bias is the subject of this paper. Selecting on it
builds the bias under test into the test sample, since the cores chosen are
those whose mis-measured masses happen to look right. The selection must
instead use quantities \emph{getsf} measures reliably and that the bias does
not touch: the background column density under the source,
${\rm PEAK^{BGF}}$; the contrast against it,
${\rm PEAK^{SBF}}/{\rm PEAK^{BGF}}$; and the measured size at half
maximum in the column-density band, AFWHM. The mass is left to follow from
whichever model those three select, and the truth table records what was
injected. The mass cannot be excluded altogether, because it is correlated
with the other three ($+0.74$ between $\log M$ and $\log\Scloud$ in Aquila),
but it enters the selection with a weight an order of magnitude below the
others: enough to keep the correlation, not enough to let the biased quantity
drive the choice.

\subsubsection{Each model must be placed at the column it was built for}
A grid model
is constructed embedded in a cloud of a particular column density $\Scloud$
and is meaningful only where the local column matches. Since the grid samples
$\Scloud$ in steps of a factor of two, the interval of local columns a model
may occupy must be one octave wide and centered on its own value,
$\Scloud/2^{1/2} \le \Sigma_{\rm local} < 2^{1/2}\,\Scloud$. A narrower
interval leaves part of the column axis belonging to no model at all: with an
interval half that width, half of the real Aquila sources, including the
source at the population's median column, fall where no model may legally sit,
and every model drawn for such a column is displaced in the same direction.

\subsubsection{Sizes and separations must refer to what is observed}
The truncation
radius $\Rbe$ is not an observable. Grid models with $\Rbe$ between $136''$
and $428''$ present a width at half maximum of only $25''$ to $50''$, because
the outer envelope is a shallow pedestal. Selecting or separating models by
$\Rbe$ therefore rejects every model above $1.8\,\msun$ while admitting nothing
that is large on the sky, and reserves exclusion zones far wider than the
models occupy. That the model width and the recovered AFWHM are the same
quantity is measured, not assumed: over $451$ injected cores matched to their
recovered counterparts the ratio has median $1.001$ and scatter $0.127$~dex,
the agreement holding to about $50''$, above which the recovered size
saturates near $50$--$60''$ through the footprint truncation of
Sect.~\ref{sec:footprint}.

\subsubsection{Contrast must be matched as the excess over the background}
The
contrast $C$ spans only $0.09$~dex over the real Aquila sources, so matching
on $\log C$ leaves it effectively unconstrained. The quantity with useful
range is $C-1 = {\rm PEAK^{SRC}}/{\rm PEAK^{BGF}}$, which spans
$0.87$~dex over the same sources.

Applied together, these give an injected population whose contrast
distribution is statistically indistinguishable from that of the real Aquila
sources (Kolmogorov--Smirnov statistic $D=0.052$ at probability $0.29$; $D$ is
the largest vertical distance between the two cumulative distributions, zero
for identical samples and unity for samples that nowhere overlap), and whose
joint distribution of column density and contrast overlaps the real one by
$0.86$ on the scale where unity is identity and zero is disjoint. Without the
four requirements the same grid and the same field give an overlap of $0.32$.
\draftnote{Values from the truth table of the population-matched run;
replace with the like-for-like comparison of extracted quantities once that
extraction is analyzed, and update Table~\ref{tab:validation} and
Fig.~\ref{fig:validation}, whose injection rows still come from the earlier,
unmatched injection of 108 sources.}


\subsection{The selection function, and what it implies for a real catalog}
\label{sec:selection}

Matching the injected population to the real one has a consequence beyond
validating the correction. Because the injected cores occupy the same region
of observable space as the real sources, and are put through the same
extraction, the fraction of them recovered measures the selection function of
that extraction where it matters, and the same selection acted on the real
sources.

Table~\ref{tab:selection} gives it. Two stages are separated: whether
\emph{getsf} detects the core at all, and whether the detection then survives
the quality criterion, QUALITY $=$ ok, which is a test on the goodness of the
spectral energy distribution fit.

\begin{table}
\centering
\caption{Fraction of injected cores recovered, against the contrast
$C = {\rm PEAK^{SBF}}/{\rm PEAK^{BGF}}$ they were injected with.
Detection means an extracted source within $30''$ of the injected position;
the quality column is the fraction of those detections that pass
QUALITY $=$ ok; completeness is the product. Pooled over the two injection
runs, $1086$ injected core instances in total.}
\label{tab:selection}
\begin{tabular}{lrrrr}
\hline\hline
contrast & injected & detected & kept of & complete- \\
         &          &          & detected & ness \\
\hline
$1.00$--$1.15$ & $156$  & $4\%$   & $14\%$ & $1\%$ \\
$1.15$--$1.30$ & $583$  & $9\%$   & $16\%$ & $1\%$ \\
$1.30$--$1.50$ & $215$  & $66\%$  & $27\%$ & $18\%$ \\
$1.50$--$2.00$ & $92$   & $98\%$  & $38\%$ & $37\%$ \\
$2.00$--$10.0$ & $40$   & $100\%$ & $20\%$ & $20\%$ \\
\hline
all            & $1086$ & $30\%$  & $27\%$ & $8\%$ \\
\hline
\end{tabular}
\end{table}

The two stages behave quite differently. Detection is nearly a step function
of contrast, rising from $9\%$ to $98\%$ between $C = 1.2$ and $C = 1.75$: a
core of low contrast is not found at all. The quality criterion removes about
three quarters of what is found, at every contrast, with only a weak trend
from $14\%$ to $38\%$ across the whole range. So low-contrast cores are lost
because they are never detected, not because their fluxes are then judged too
uncertain.

The detection threshold responsible is deliberate, and the same injections
show what it buys. In a simulated cloud containing no sources at all, a low
threshold extracts large numbers of background fluctuations as if they were
sources. The threshold is set to suppress them, on the principle that a
smaller sample of well-measured sources is more useful than a larger one
diluted by peaks whose fluxes carry uncertainties comparable with their
values. In the injected fields analysed here the price is small and the
benefit large: of the $357$ sources extracted, $151$ correspond to injected
cores and $203$ to real Aquila sources, while only $3$, under $1\%$,
correspond to neither and are candidates for spurious detections. The
incompleteness in Table~\ref{tab:selection} is therefore the deliberate cost
of that purity, not a defect, and the injection machinery described here is
also the tool with which the trade-off could be re-optimised, by repeating the
measurement at a lower threshold and recording both numbers.

Two remarks on the comparison. The contrast against which the completeness is
tabulated is the one each core was injected with,
$1 + {\rm ICSDbs}/\Sigma_{\rm local}$, whereas a real source is
described by the contrast \emph{getsf} reports. Over the $328$ detected
injected cores the two agree to a median factor of $1.09$ in $C-1$, with a
scatter of $0.133$~dex, so the two axes may be compared directly at this
level.

The relevance to a real catalog follows from where the real sources sit. Over
the $1030$ Aquila sources the median contrast is $1.235$, which is inside the
band detected at $9\%$. The catalog is therefore highly incomplete precisely
where the bulk of its own sources lie, and the cores it is missing are the
low-contrast ones.  Whether those missing objects are prestellar cores at all
is a separate question: a low-contrast peak in a real cloud may be a genuine
density enhancement without being a core of the kind this work addresses.
What the injections establish is only that objects which are such cores,
at that contrast, would not have been found. Since contrast correlates with background column at
$+0.49$ in the logarithm, and reported mass with column at $+0.74$, the
missing population is preferentially of low mass on faint background.

This interacts with the correction in an unhelpful direction.
$\CM$ is largest for the sources whose footprints truncate the most flux,
which are the extended, low-contrast ones, and those are exactly the cores the
selection removes. A corrected real catalog therefore still understates the
total mass in cores, by more than the correction itself would suggest: the
correction repairs the masses of the cores that were found, and says nothing
about those that were not.

We stop short of inverting the selection to recover the underlying
population. Doing so is not a matter of reweighting each catalog source by the
inverse of its completeness, for two reasons. The completeness is measured
against the contrast a core was injected with, whereas a real source is known
only through the contrast the extraction reports for it, which the same
selection distorts. And the injected population was matched to the real
detected catalog, which is itself the survivor of this selection, so
using it to undo that selection is circular. The consistent procedure is
forward modeling: to postulate an underlying population, inject it, apply the
measured selection, and compare the predicted detections with the observed
catalog, iterating on the postulate. That is a larger undertaking than this
paper, though it needs no machinery beyond what is described here.
\draftnote{Numbers from the two runs of 2026 August 10. Regenerate when the
further injection realizations are available; the two lowest contrast bins
rest on $7$ and $50$ detections respectively.}


\subsection{Why the three-observable correction did not reduce scatter,
and what the fourth observable changes}
\label{sec:scatterfloor}

Table~\ref{tab:validation} shows a clear split: the three-observable
corrector's scatter is worse than not correcting at all in every test, while
the four-observable corrector's is better in two of three and closely
matched in the third. Both observations have the same explanation. A
corrector predicts a single number, $\CM$, from a fixed set of
observables. Two sources with identical values on every observable used
can still have different individual mass biases, because those observables
do not capture every detail relevant to the bias; the correction can only
remove the part that is predictable from what it is given, not the part
that is not. This has a clean statistical signature: multiplying an
already-noisy measurement by a correction factor of finite precision
cannot decrease the total variance, and increases it by the corrector's
own prediction variance if that variance is independent of the original
measurement noise, since variances of independent quantities add in log
space,
\begin{equation}
  \sigma_{\rm corrected}^2 \approx \sigma_{\rm uncorrected}^2
    + \sigma_{\CM}^2 .
  \label{eq:scatterbudget}
\end{equation}
For the three-observable corrector, Sect.~\ref{sec:fourthobs} showed
directly why $\sigma_{\CM}$ is large: $(\Scloud,\zeta,
\phi)$ leave a genuine, non-shrinking degeneracy in $\CM$ of a
factor of $1.1$--$1.2$ even between arbitrarily close neighbors
(Fig.~\ref{fig:degeneracy}), because they carry no information about a
source's absolute physical scale. Solving Eq.~(\ref{eq:scatterbudget}) for
the three-observable corrector's own scatter using the measured
before/after values gives a factor of about $1.3$ in the injection test
and the more demanding node-\emph{loo} test, and a smaller factor of about
$1.1$ in the subdivision test, which spans a narrower range of the grid's
diversity per node than the full production grid does. Adding the fourth
observable does not eliminate $\sigma_{\CM}$, since interpolation error
and real measurement noise remain, but it removes the dominant, avoidable
contribution, which is why Table~\ref{tab:validation}'s four-observable
column shows scatter below the uncorrected baseline rather than above it.

Whether the residual scatter matters depends on what the corrected masses
are used for. Write $\Mref$ for the known reference mass a test
compares against ($\Mref$ for the injection test, $\Mbe$ for
the subdivision and \emph{loo} tests) and $\delta \equiv \log_{10}(M/\Mref)$
for the log-ratio of a single sample's mass (uncorrected or corrected, as
context requires) to that reference, and
$\delta_{\rm rms}\equiv\langle\delta^2\rangle^{1/2}$ for its root-mean-square
value over a sample: because $\delta=0$ is exact agreement with the
reference mass, $\delta_{\rm rms}$ combines bias and scatter into the
total error against the truth, unlike the standard deviation of $\delta$,
which measures only the spread about the sample's own mean and says
nothing about how far that mean sits from zero. Even with the fourth
observable, a single source's point-estimate precision is not
guaranteed to improve on simply knowing the population's average bias
(Table~\ref{tab:validation} and Fig.~\ref{fig:validation} report this
honestly rather than optimizing the figure of merit that flatters the
method). For a population statistic computed over many sources -- a total
cloud mass, a mean mass in some environmental bin, the normalization or
shape of a mass function -- the distinction matters a great deal: a
systematic bias, uncorrected, biases every such statistic by the same
factor regardless of sample size, while the correction's added per-source
scatter is close to symmetric in log space and mass-independent
(Fig.~\ref{fig:validation}a, c, e), so it averages down with sample size
in linear statistics and contributes only second-order effects to the
shape of a fitted power law. Removing the
former at a small mildly negative cost to the latter is the intended
trade, and the reason the correction is built and validated as a bias
correction rather than as a per-source precision improvement.

\subsection{Which factor drives the scatter}
\label{sec:factorscatter}

Because $\CM=\Fflux\times \Fbg\times \Ftemp$
(Sect.~\ref{sec:footprint}) is built from two independently interpolated
factors, $\Fflux$ and $\Fbg$, with $\Ftemp$ obtained as
the residual $\CM/(\Fflux\Fbg)$, the scatter budget of
Sect.~\ref{sec:scatterfloor} can be attributed to a specific mechanism. We
repeat the single-cross-section subdivision test that preceded the
current one (Sect.~\ref{sec:validation}) for each factor separately,
comparing its interpolated prediction against its true value at each
withheld node (Fig.~\ref{fig:factordecomp}).

\begin{figure*}
  \centering
  \includegraphics[width=0.72\linewidth]{fig_factor_decomposition.pdf}
  \caption{Prediction accuracy of the three factors composing $\CM$, on
  the three-observable corrector, using the original ($n=873$),
  single-cross-section subdivision test that preceded the fourth
  observable and motivated its adoption. (\emph{top}) flux recovery,
  $\Fflux$. (\emph{middle}) background subtraction, $\Fbg$.
  (\emph{bottom}) temperature-gradient bias, $\Ftemp$. (\emph{left})
  predicted-over-true ratio versus $\Mbe$. (\emph{right}) the
  accuracy CDF of each factor, with the scatter of its predicted-over-true
  ratio given in the title as a multiplicative factor.}
  \label{fig:factordecomp}
\end{figure*}

The three factors are not equally well determined. $\Ftemp$ is
predicted almost exactly (within a factor of $1.05$): the internal
temperature-gradient bias is a smooth, robust function of
$(\Scloud,\zeta,\phi)$ with essentially no
residual scatter beyond interpolation noise. $\Fbg$ is
intermediate (within a factor of $1.30$). $\Fflux$ is by far the
least well determined (within a factor of $1.65$), consistent with it
depending on the finer structure of the per-band footprint
(Sect.~\ref{sec:footprint}), which is more sensitive to a source's
specific profile shape than the other two factors. Because it is the
logarithm of each factor that behaves as an independent Gaussian-like
error, the corresponding log-space variances (not the factors
themselves) are what combine in quadrature, and doing so reproduces
$\CM$'s own scatter to within $0.4\%$, confirming the three factors are
being fitted independently rather than trading error against each other.
In this decomposition, $\Fflux$ alone accounts for $77\%$ of
$\CM$'s prediction variance, $\Fbg$ for $21\%$, and $\Ftemp$
for under $1\%$.

$\Fflux$ being the dominant source of scatter does not mean it
should be dropped from the correction. Its own median predicted-over-true
ratio is $0.996$, indistinguishable from unbiased, and it carries the
single largest median contribution to $\CM$ among the three factors
(Sect.~\ref{sec:correction}); removing it would reintroduce the flux-loss
bias the footprint model of Sect.~\ref{sec:footprint} was built to
correct, trading a small scatter reduction for a large, systematic
mass under-estimate. This decomposition, run on the three-observable
corrector, is what motivated the search for a fourth observable
(Sect.~\ref{sec:fourthobs}): a quantity that reduces $\Fflux$'s
scatter propagates directly into $\CM$'s scatter, while equivalent effort
on $\Ftemp$ would not, and the angular size adopted in
Sect.~\ref{sec:fourthobs} does exactly this, since the flux-recovery
physics depends on the source's scale through the same self-shielding
column that sets its temperature structure.

\subsection{An output cap was tested and superseded}
\label{sec:outputcap}

Before adopting the fourth observable, we tested a less invasive
alternative: capping the corrected mass from above, since the
three-observable corrector's excess scatter (Sect.~\ref{sec:scatterfloor})
is dominated by occasional large over-corrections rather than a uniformly
broadened distribution. On the validation samples available at the time
(the earlier, single-cross-section subdivision test and the pre-fix
production grid of Sect.~\ref{sec:interp}), capping
$\Mcorr/\Mref$ at a factor of $1.5$ recovered much of the
lost accuracy (the fraction within a factor of two rose to
$92$--$96\%$ across the three tests, with the median unaffected), but this
is a diagnostic result, not a production fix: the cap was applied to the
ratio against a known truth, available only in a validation setting, and
translating it to a cap on $\CM$ itself for real sources requires choosing
and separately validating a threshold from the $\CM$ distribution alone.
The fourth observable achieves a comparable or
better result (Table~\ref{tab:validation}) without this extra step, by
supplying the missing information directly rather than truncating its
symptoms, and is the method we adopt.

% =====================================================================
\section{Discussion}
\label{sec:discussion}

\subsection{Resolution dependence of the surface-density band}
\label{sec:resolution}

The grid, and hence the correction, is calibrated at the $13.5''$
resolution of the high-resolution surface-density image used throughout
this work. Applying the correction to an extraction that instead relies on
an $18.2''$-resolution surface-density proxy, without rebuilding the grid at
that resolution, biases two of the three observables low, and the bias is
largest for the most compact cores: convolving the grid's surface-density
profiles from $13.5''$ to $18.2''$ and recomputing $\zeta$ and
$\phi$ axes at fixed truncation radius gives a bias of $-33\%$
($\zeta$) and $-18\%$ ($\phi$) for cores smaller than the
$18.2''$ beam, falling to about $-1\%$ for cores larger than $\sim\!80''$
(Appendix~\ref{app:resolution}). The method is not intrinsically tied to
$13.5''$: applying it at coarser resolution requires re-tabulating the
model recovery grid with the surface-density beam set to that resolution
(re-convolving the existing radiative-transfer model images, with no new radiative-transfer
calculation), not substituting a coarser-resolution catalog into the
present, finely-calibrated tables.

\subsection{The low-mass edge and the correction's domain}
\label{sec:lowmass}

The correction is not applicable at arbitrarily low mass. Below a
column-dependent limit, the coldest, faintest cores lose enough flux that
the recovered observables no longer determine $\Mbe$ uniquely, and
the correction becomes degenerate; this limit rises steeply with cloud
background column, from a few hundredths of a solar mass in the lowest-
column part of the grid to several solar masses in the highest-column part.
The correction is applied only above this limit and is not extrapolated
below it.

\subsection{What the method does, and does not, correct}
\label{sec:scope}

The correction targets a source that is a genuine, single, pressure-
truncated BE sphere observed with the stated systematics. It does not
correct for unresolved multiplicity: at coarse resolution and large
distance, several genuine cores may be blended into one \emph{getsf}
detection, and a large correction factor applied to such a blend returns the
mass of the equivalent single object that would reproduce the blend's
observables, not the mass of any individual member. Distinguishing a
genuinely massive single core from an unresolved multiple is outside the
scope of this method and requires higher-resolution follow-up.

% =====================================================================
\section{Conclusions}
\label{sec:conclusions}

We have described a radiative-transfer-calibrated correction for the two
dominant biases in \emph{Herschel} SED-derived prestellar core masses -- internal
temperature-gradient bias and footprint-limited flux loss -- built on a
grid of pressure-truncated Bonnor--Ebert models processed exactly as a real
\emph{getsf} extraction. A first, three-observable version of the
correction, using a self-consistent per-band detection footprint set by a
constant source signal-to-noise ratio and the beam-convolved power
spectrum of the cloud's structural noise, removes the systematic bias but
leaves a genuine degeneracy: sources sharing all three observables can
still require correction factors differing by a factor of two or more,
because the observables are constructed to be scale-invariant while the
correction itself is not. We identified this degeneracy directly, by
showing that the spread in the true correction factor between
near-identical training samples does not shrink as the samples are made
arbitrarily close, and traced it to the missing scale of the source.
Adding the angular truncation radius as a fourth observable, matched on the
real-source side by the measured footprint radius in the same angular
units, removes most of this degeneracy. Because the correction factor is a
ratio of two masses that carry the same power of distance, it is
distance-free by construction, and the fourth axis is used as an angle;
converting it to a physical length at each source's own distance was found
to introduce a spurious trend of the correction with distance.

Validated on three independent tests -- synthetic cores carried through
the full injection, extraction, and correction chain; an independent
model grid spanning the full range of core concentration rather than a
single stability class; and a leave-one-out test repeated at every node of
the production grid -- the four-observable correction is unbiased to
within a few per cent, with $85$--$98\%$ of sources recovered within a
factor of two of the reference mass, and its scatter is at or below the
uncorrected level in every test, a property the three-observable version
did not have. Application of the method to observed cloud catalogs, and
its consequence for the derived core mass function, are presented in a
companion paper.

\begin{acknowledgements}
\draftnote{Acknowledgements to be added.}
\end{acknowledgements}

\bibliographystyle{aa}
\bibliography{rt_mass_correction}

\appendix

% =====================================================================
\section{Mass conservation in the grid construction}
\label{app:massconserv}

The radial grid used for the radiative-transfer calculation is defined once
per model and written to the radiative-transfer code's grid file; every subroutine that
needs the grid reads this file rather than re-deriving cell boundaries
independently, which was found to be the root cause of a boundary-mass
error in an earlier version of the pipeline, where a stale, independently
computed duplicate of the grid drifted out of synchronization with the one
actually used for the transfer calculation. The BE truncation radius
$R_{\rm BE}$ is inserted as an explicit wall in the radial grid, rather than
approximated by snapping to the nearest pre-existing cell boundary; this
preserves the grid's logarithmic spacing away from the truncation radius
while ensuring that no single cell straddles the BE/cloud interface. Cells
are then unambiguously assigned to the BE profile (inner wall radius below
$R_{\rm BE}$) or the uniform cloud (inner wall radius at or above
$R_{\rm BE}$), a partition that is complementary and leaves neither a
vacuum gap nor a double-counted shell at the boundary. With this
construction, the mass obtained by summing the density in every grid cell
of a model agrees with the analytic BE mass integral to better than
$0.1\%$ for every model in the grid.

% =====================================================================
\section{Why the correction is not indexed on reported mass}
\label{app:distance}

An earlier formulation of the correction matched a source to the grid using
its reported mass as one of the matching variables, together with the
background column and $\zeta$. Because the reported mass of a source
of fixed physical structure scales approximately as the square of distance,
this formulation matches sources at different distances to systematically
different grid models even when their intrinsic structure is identical:
the correction factor derived for the same source observed at $260\,$pc and
at $1700\,$pc differs by an amount driven purely by distance, not by
physics. Removing the mass axis and matching instead on the three
distance-invariant observables of Sect.~\ref{sec:observables} removes this
failure mode entirely, with no loss of accuracy: a node-level leave-one-out
test on the recovery tables gives the same median error for the two
formulations to within a few per cent, while only the distance-invariant
formulation gives a correction factor that is stable across clouds spanning
a wide range of distance (demonstrated for seven clouds in the companion
application paper).

The fourth observable of Sect.~\ref{sec:fourthobs} reintroduces a distance
dependence, but not the one rejected here. Reported mass scales as
distance squared and is itself the biased quantity the correction exists
to fix, so using it as a matching variable means comparing a source
against the grid using the very number that is wrong. The angular
footprint radius adopted instead carries no distance dependence at all,
and it is not the quantity under correction: it is not biased by the
mechanisms of Sect.~\ref{sec:footprint} in the way the reported mass is.

The objection to the reported mass is therefore narrower than it once
appeared, and it is worth stating what it does and does not cover. Under
the leave-one-out test of Sect.~\ref{sec:validation}, with each candidate
used consistently for training and query, the reported mass is in fact the
better fourth axis of the two, giving a residual in $\log_{10}\CM$ of
$0.098$~dex against $0.118$ for the footprint radius. What disqualifies it
is not its performance but its failure mode: it inherits the very bias the
correction exists to remove, so an error in it displaces the source along
the corrector's own axis, whereas an error in an angular size cannot,
since angles are measured directly. \draftnote{If a future version accepts a
mass-bearing axis, this appendix needs rewriting rather than deleting: the
performance claim and the failure-mode claim point in opposite
directions.}

% =====================================================================
\section{Interpolation accuracy of the model grid}
\label{app:interp}

The nominal factor-of-two node spacing of the production grid is adequate
where the tabulated observables vary smoothly with the grid parameters, but
can be too coarse where an observable changes character -- for example
where the $70\,\mu$m signed center intensity of a cold, embedded core
crosses from emission to absorption against the background, or where a
photometric band transitions from optically thin to optically thick.
Interpolation accuracy is verified by a hold-out test: the interpolated
observable vector at a withheld grid node is compared against a dedicated
radiative-transfer run at that exact node, with particular attention to the
nodes that bracket the transitions described above, and the local grid
spacing is refined wherever any observable's interpolation error exceeds a
few per cent.

% =====================================================================
\section{Resolution-dependence estimate}
\label{app:resolution}

The bias quoted in Sect.~\ref{sec:resolution} is estimated by convolving
each grid node's surface-density-band radial profile from the production
$13.5''$ resolution to $18.2''$ with a one-dimensional Gaussian kernel of
width $(18.2''^2-13.5''^2)^{1/2}/2.3548$, and recomputing $\zeta$
and $\phi$ observables from the convolved profile at the same
truncation radius as the true $13.5''$ value. This approximates the true
two-dimensional convolution of the model image by a one-dimensional
convolution of its azimuthally-averaged radial profile, which is adequate
to establish the sign and approximate magnitude of the resolution bias but
is not a substitute for re-tabulating the recovery grid with the
surface-density band convolved to $18.2''$ in two dimensions, which is the
recommended procedure for applying the method at that resolution.

\section{Constructing a population-matched injection test}
\label{app:injection}

This appendix gives the procedure summarized in Sect.~\ref{sec:injpop}, in the
order it is carried out. The code is available with the paper; what follows is
the reasoning behind each step and the numbers needed to judge it.

\subsection{The target table}
One row per real source, from the \emph{getsf}
extraction of the field to be imitated, carrying the reported SED mass, the
background column ${\rm PEAK^{BGF}}$, the contrast
${\rm PEAK^{SBF}}/{\rm PEAK^{BGF}}$, and the width at half maximum in
the $13.5''$ column-density band. For Aquila this gives $1030$ sources.

The table may be drawn from a wider field than the one the cores are injected
into. The injections here are made in a sub-field that is denser than Aquila
as a whole, with a median source background of
$1.0\times10^{22}\,{\rm cm^{-2}}$ against $5.0\times10^{21}\,
{\rm cm^{-2}}$ for the full field. Split by background column, however, the
two populations have the same median contrast and the same median reported
mass in every octave; only the relative number of sources at each column
differs, and that is set by the map the cores are placed in, not by the table.
Using the full field therefore doubles the statistics without importing a
bias. Taking the marginal distribution from the wrong field, by contrast, is a
genuine error, and was why an earlier version of this test appeared to place
cores at systematically too high a column.

\subsection{Choosing a model for each core}
A reported mass $M$ is drawn from
${\rm d}N/{\rm d}M \propto M^{-\alpha}$; one of the $40$ real sources
nearest to it in $\log M$ is picked at random, and its background column
$\Sigma$, contrast $C$ and size $A$ become this core's target. Taking all
three from a single real source preserves their mutual correlations, which in
Aquila are strong ($+0.74$ between $\log M$ and $\log\Sigma$, $+0.49$ between
$\log\Sigma$ and $\log(C-1)$); drawing them independently would impose the
population-median contrast at every column, which is wrong at both ends.

Among the grid nodes eligible at that column, and whose width at half maximum
lies in the permitted range, the node is chosen by minimizing
\begin{equation}
  d = w_{M}\left[\log\frac{M}{\Msed}\right]^{2}
    + w_{C}\left[\log\frac{C-1}{C_{n}-1}\right]^{2}
    + w_{A}\left[\log\frac{A}{{\rm FWHM}_{n}}\right]^{2},
  \label{eq:injscore}
\end{equation}
with $\Msed$ the node's own reported four-band SED mass,
${\rm FWHM}_{n}$ its width at half maximum, and
$C_{n} = 1 + {\rm ICSDbs}/\Sigma$ the contrast it would show at the column
it is about to occupy, ${\rm ICSDbs}$ being its background-subtracted
central column. The choice is not the strict minimum: a node is drawn at
random among those of nearly equal $d$, with probability falling off as $d$
worsens, so that repeated draws of similar targets do not collapse onto one
node.

Two details matter. The mass is compared with the node's own reported SED mass
rather than with its true mass $\Mbe$: over the eligible nodes the
ratio $\Mbe/\Msed$ has median $1.105$, so comparing with
the truth would build in a systematic offset of $0.043$~dex toward models that
are too light. And the weights are unequal: with $w_{C}=w_{A}=1$, raising
$w_{M}$ from $0.1$ to $1$ degrades the contrast agreement from $D=0.063$ at
probability $0.11$ to $D=0.174$ at probability $10^{-10}$, and the overlap in
the plane of column and contrast from $0.89$ to $0.85$.

\subsection{Choosing a position for each core}
A position is admissible if the model stamp fits
inside the coverage, if the local column lies in the model's octave, and if
the center clears the footprints of the real sources by one $500\,\mu$m beam
radius. The margin applies to the center alone. Requiring the whole model to
clear them, which is what taking the model radius to be $\Rbe$ implies, turns
footprints covering $6.5\%$ of the field into an exclusion zone covering most
of it. What must be prevented is an injected core being centered on a real
one, making the extraction a blend; overlap of the outer envelope with a real
footprint is harmless and in a real cloud unavoidable.

Each core is steered to the column drawn for it individually, by ordering the
admissible positions by $|\log(\Sigma_{\rm local}/\Sigma)|$ and taking
the first that satisfies the separation criterion. Steering all cores of one
node to the median of their drawn columns instead, the obvious economy since
the admissible set depends only on the node, collapses the spread of columns
within each node: with about eight cores per node the rank correlation between
contrast and column falls to $0.26$ against $0.51$ in the real sample. Handled
individually, the drawn column is realized exactly, the median ratio of
realized to drawn being $1.000$.

Two injected cores of the same size class may not lie closer than
$1.3\,(r_{i}+r_{j})$, with
$r=\left(({\rm FWHM}/2)^{2}+\theta_{500}^{2}\right)^{1/2}$. The criterion binds within
a class only, because a large faint model is barely detectable against the
cloud and does not disturb a small core lying on it, whereas two overlapping
extended models would double the column over an extended area.

\subsection{The residual mismatch in size}
The measured size remains discrepant, and the
cause is the coverage of the model grid rather than the placement. Once mass
and contrast are fixed a Bonnor--Ebert sphere has no freedom left in size: a
sphere of $0.1\,\msun$ showing a contrast of $1.23$ against
$5\times10^{21}\,{\rm cm^{-2}}$ must be large and diffuse, whereas real
Aquila sources of that mass and contrast are measured at $31$--$38''$. At the
embedding column $6\times10^{21}\,{\rm cm^{-2}}$ the faintest grid model
between $26''$ and $35''$ has contrast $1.87$, and between $45''$ and $55''$
contrast $2.71$, both far above what the real sources need; only the size
classes near $42''$ and near $70''$ hold models faint enough, and the injected
cores accordingly pile up there and are absent below $26''$, where a quarter
of the real sources lie.

Two checks confirm the diagnosis. Restricting the models to $15$--$55''$, to
remove the excess of large injected cores, makes the agreement worse, not
better ($D=0.052\to0.109$ in contrast, overlap $0.86\to0.82$): the models
between $55''$ and $80''$ are needed not because real sources are that large
but because on this grid they are the only ones faint enough. And correcting
the injected sizes to what the extraction would report changes the size
statistic only from $D=0.425$ to $0.420$, so the mismatch is not an artifact
of comparing a model parameter with a measured quantity. The grid is missing
faint compact models: nodes of low central column,
${\rm ICSDbs}\sim1$--$2\times10^{21}\,{\rm cm^{-2}}$, below about
$35''$, and faint models at all sizes in the lowest embedding rung.
\draftnote{Grid extension planned; re-run and update.}

\subsection{The realized mass function}
The mass function of the cores actually
injected is not the one drawn. With $\alpha=2$, so
${\rm d}N/{\rm d}\log M \propto M^{-1}$, the true masses assigned follow
$M^{-1.47\pm0.14}$ above $0.20\,\msun$, the steepening arising at the
high-mass end where the grid has no node of the required mass at a realistic
size and contrast. The input slope can be adjusted to compensate, at no cost
in the per-core matching, but at a real cost in the population match: the
shallower input draws more massive real sources as targets, and these sit on
denser, higher-contrast backgrounds, so the contrast statistic rises from
$D=0.052$ to $0.228$. The two goals conflict because mass and column are
strongly correlated in the real population. We keep the population match and
take the realized mass function from the truth table; no step of the analysis
uses the nominal slope.


\end{document}
